\documentclass[11pt,letterpaper]{amsart}
\usepackage[margin=2.5cm]{geometry}
\usepackage{setspace}
\onehalfspacing
\raggedbottom
%---------------------------------------------
%\usepackage{backref} % The package hyperref was calling a command from backref, but backref was not added before, and so the compiler kept producing an error. Including the package fixed the error and *hopefully* did not mess anything up.

%\usepackage{graphicx}
\usepackage{amsmath}
\usepackage{mathtools}
\usepackage{amssymb}
%\usepackage{amsfonts}
\usepackage{amsthm}
\usepackage{todonotes}
\usepackage{bbm}
\usepackage{stmaryrd}
%\usepackage{cancel}
\usepackage[normalem]{ulem}
%\usepackage[all]{xy}
%\usepackage[neveradjust]{paralist}
%\usepackage{enumitem}
%\usepackage{multicol}
%\usepackage{mathlocrsfs}
\usepackage{mathdots}
\usepackage{float}
\usepackage[all]{xy}

%\usepackage{enumitem}

\usepackage[pagebackref,colorlinks=true,linkcolor=cyan,citecolor=magenta]{hyperref}

\usepackage[alphabetic,nobysame, initials, backrefs]{amsrefs}
\usepackage{tikz-cd}
\usepackage{tikz}
\usepackage{nicematrix}
\usepackage{xcolor}
%---------------------------------------------------------------
% Revision markup: everything added or altered in this revision is
% printed in red via \rev{...}; \revnote{...} flags a point that still
% needs an author decision.  Delete these two macros' colouring
% (i.e. \newcommand{\rev}[1]{#1}) to produce the clean version.
%---------------------------------------------------------------
% The original \todo notes from newversion1.tex are retained, each with a
% [disposition:] line saying how it was handled.  Switch them all off with
%     \usepackage[disable]{todonotes}
% in the \usepackage line above.
\newcommand{\rev}[1]{{\color{red}#1}}
\newcommand{\revnote}[1]{{\color{red}[\textsf{\textbf{note:}} #1]}}
% \old{...} reproduces, in blue, a statement from newversion1.tex that has been
% replaced because it was not correct as written; delete these blocks once checked.
\newcommand{\old}[1]{{\color{blue}#1}}
%\usepackage{blkarray}
%\usepackage{scalerel,stackengine}

% \UseRawInputEncoding

% \stackMath
% \newcommand\reallywidehat[1]{%
% \savestack{\tmpbox}{\stretchto{%
%   \scaleto{%
%     \scalerel*[\widthof{\ensuremath{#1}}]{\kern-.6pt\bigwedge\kern-.6pt}%
%     {\rule[-\textheight/2]{1ex}{\textheight}}%WIDTH-LIMITED BIG WEDGE
%   }{\textheight}% 
% }{0.5ex}}%
% \stackon[1pt]{#1}{\tmpbox}%
% }



% \setlength{\hoffset}{-0.35in} \addtolength{\textwidth}{0.70in}
% \setlength{\voffset}{-0.50in} \addtolength{\textheight}{1.00in}


\tikzset{
  symbol/.style={
    draw=none,
    every to/.append style={
      edge node={node [sloped, allow upside down, auto=false]{$#1$}}}
  }
} 

\usepackage{imakeidx}
\makeindex[columns=2]
\newcommand{\BH}[1]{{\hyperpage{#1}}\normalsize}
\newcommand{\IN}[1]{\index{#1|BH}}

\newcommand{\IC}{\mathcal{IC}}
\newcommand{\Loc}{\operatorname{Loc}}
%\usepackage{color}
\newcommand{\ceq}{\, := \, }
\newcommand{\tq}{\ \vert\ }

\newcommand{\N}{\mathbb{N}}
\newcommand{\Z}{\mathbb{Z}}
\newcommand{\Q}{\mathbb{Q}}
\newcommand{\R}{\mathbb{R}}
\newcommand{\BC}{\mathbb{C}}

\newcommand{\CP}{\mathbb{CP}}
\newcommand{\Cat}{\mathcal{C}}
\renewcommand{\P}{\mathbb{P}}
\renewcommand{\H}{\mathbb{H}}
\newcommand{\D}{\mathbb{D}}
\newcommand{\A}{{\mathbb A}}

\newcommand{\ord}{{\mathrm{ord}}}

\newcommand{\p}{\mathfrak{p}}
%\newcommand{\q}{\mathfrak{q}}
\renewcommand{\a}{\mathfrak{a}}
\newcommand{\m}{\mathfrak{m}}


\newcommand{\Acal}{\mathcal{A}}
\newcommand{\Tcal}{\mathcal{T}}
\newcommand{\Ecal}{{\mathcal E}}
\newcommand{\Fcal}{{\mathcal F}}
\newcommand{\Homcal}{{\mathcal {H}om}}
\newcommand{\Gcal}{{\mathcal G}}
\newcommand{\Hcal}{{\mathcal H}}
\newcommand{\Ocal}{{\mathcal O}}

\newcommand{\Tau}{{\mathcal T}}
\newcommand{\B}{{\mathcal B}}
\newcommand{\OO}{{\mathcal O}}

\newcommand{\inv}{^{-1}}
\newcommand{\sm}{\setminus}
\newcommand{\ol}{\overline}

\newcommand{\SL}{\mathrm{SL}}
\newcommand{\GL}{\mathrm{GL}}
\newcommand{\SO}{\mathrm{SO}}
\newcommand{\Sp}{\mathrm{Sp}}
\newcommand{\St}{\mathrm{St}}
\newcommand{\dualgroup}[1]{{#1}^\vee}
\newcommand{\irrepdual}[1]{\widehat{#1}}
\newcommand{\Lgroup}[1]{{\hskip-2 pt \,^L\hskip-1pt{#1}}}
\newcommand{\Der}{\mathrm{Der}}
\newcommand{\reg}{\mathrm{reg}}
\newcommand{\gen}{\mathrm{gen}}
\newcommand{\NEvs}{\mathrm{NEvs}}
\newcommand{\Evs}{\mathrm{Evs}}
\newcommand{\pure}{\mathrm{pure}}
\newcommand{\st}{\mathrm{st}}
\newcommand{\std}{\mathrm{std}}
\newcommand{\pr}{\mathrm{pr}}
\newcommand{\QST}{\mathrm{QST}}
\newcommand{\sgn}{\mathrm{sgn}}


% Desiderata
\newcommand{\Sta}{\textbf{(St)}}
\newcommand{\ECILS}{\textbf{(ECI}$^{\mathrm{LS}}$\textbf{)}}
\newcommand{\Funtheta}{\textbf{(Fun$_{\theta}$)}}
\newcommand{\LC}{\textbf{(LC)}}
\newcommand{\Sha}{\textbf{(Sha)}}
\newcommand{\ECIKS}{\textbf{(ECI}$^{\mathrm{KS}}$\textbf{)}}
\newcommand{\ECI}{\textbf{(ECI)}}
\newcommand{\ECR}{\textbf{(ECR)}}
\newcommand{\Rel}{\textbf{(Rel)}}


% Constant sheaf


%\newcommand{\1}{{\mathbb{1}}}

\newcommand{\ul}[1]{\underline{#1}}


\newcommand{\mfr}[1]{\mathfrak{#1}}

\newcommand{\Fr}{\mathrm{Fr}}
\newcommand{\RG}{\mathrm{G}}
\newcommand{\RH}{\mathrm{H}}
\newcommand{\Ind}{\mathrm{Ind}}
\newcommand{\Res}{\mathrm{Res}}
\newcommand{\LLC}{\mathrm{LLC}}
\newcommand{\Per}{\mathrm{Per}}
\newcommand{\Shv}{\mathrm{Shv}}
\newcommand{\mic}{\mathrm{mic}}
\newcommand{\codim}{\mathrm{codim}~}
\newcommand{\temp}{\mathrm{temp}}


\newcommand{\Trans}{\mathrm{Trans}}
\newcommand{\KS}{\mathrm{KS}}
\newcommand{\Lift}{\mathrm{Lift}}

% \newcommand{\Irr}{\mathrm{Irr}}

\newcommand{\ABV}{{\mbox{\raisebox{0pt}{\scalebox{0.6}{$\mathrm{ABV}$}}}}}

\newcommand{\Art}{{\mbox{\raisebox{1pt}{\scalebox{0.6}{$\mathrm{Art}$}}}}}

\newcommand{\qs}{{\mbox{\raisebox{1pt}{\scalebox{0.6}{$\mathrm{,\circ}$}}}}}

\newcommand{\etaABV}[2]{{\eta^{\ABV #2}_{#1}}}


\newcommand{\etaArt}[2]{{\eta^{\Art #2}_{#1}}}

% Sheaf Macros
\newcommand{\locsys}[2]{\underline{#1}_{#2}}
\newcommand{\stsh}[2]{{#1}_{#2}}




% \newcommand{\Irr}[2]{\IfNoValueTF{#2}{\mathrm{Irr}(#1)}{\mathrm{Irr}(#1,#2)}}

 \newcommand{\Irr}[1]{\mathrm{Irr}({#1})}


% \usepackage{xparse}

% \NewDocumentCommand{\Irr}{mm}{%
%   \IfNoValueTF{#2}{%
%     \mathrm{Irr}(#1)%
%   }{%
%     \mathrm{Irr}(#1,#2)%
%   }%
% }




\newcommand{\STD}[2]{{\mu_{#1}^{#2}}}
\newcommand{\Rphi}[1]{\mathrm{R}\Phi_{#1}}


\newcommand{\SR}{{\Delta}}
\newcommand{\sr}{{\Delta}}


%\newcommand{\pureABV}[2]{{\eta^{\ABV, \mathrm{pure}}_{#1,#2}}}
%\newcommand{\pureArt}[2]{{\eta^{\Art, \mathrm{pure}}_{#1,#2}}}


\newcommand{\Transfer}[2]{{\Trans^{#2}_{#1}}}
\newcommand{\tTransfer}[2]{\widetilde{\Trans}^{#2}_{#1}}

\newcommand{\Lyft}[2]{{\Lift^{#2}_{#1}}}
\newcommand{\Uber}[2]{{\mathrm{Uber}}_{#1}^{#2}}


\renewcommand{\sp}{\text{sp}}

\newcommand{\PRep}{\underline{\mathrm{PRep}}}
\newcommand{\Rep}{\mathrm{Rep}}
\newcommand{\Irrep}{\mathrm{Irrep}}
\newcommand{\Seg}{\underline{\mathrm{Segments}}}

\newcommand{\half}[1]{\frac{#1}{2}}

\renewcommand{\implies}{\Rightarrow}
\newcommand{\implied}{\Leftarrow}

\newcommand{\dirlim}{\varinjlim}
\newcommand{\invlim}{\varprojlim}

\newcommand{\comment}[1]{}
\newcommand{\EE}{\mathcal{E}}
\newcommand{\FF}{\mathcal{F}}

\newcommand{\CC}{\mathbb{C}}

\newtheorem{thm}{Theorem}[section]
\newtheorem{cor}[thm]{Corollary}
\newtheorem{lemma}[thm]{Lemma}
\newtheorem{prop}[thm]{Proposition}
\newtheorem {conj}[thm]{Conjecture}
\newtheorem {assu}[thm]{Assumption}
\newtheorem {ques/conj}[thm]{Question/Conjecture}
\newtheorem {hypothesis}[thm]{Hypothesis}
\newtheorem {ques}[thm]{Question}
\newtheorem{defn}[thm]{Definition}
\newtheorem{remark}[thm]{Remark}
\newtheorem{prpt}[thm]{Property}
\newtheorem{exmp}[thm]{Example}
\newtheorem{fact}[thm]{Fact}
\newtheorem{algo}[thm]{Algorithm}
\newtheorem{conv}[thm]{Convention}
\newtheorem{problems}[thm]{Problems}
\newtheorem{appl}[thm]{Application}
\newtheorem{assumption}[thm]{Working Hypotheses}

\newtheorem{desideratum}[thm]{Desideratum}

\newtheorem*{globalconj*}{Global Conjecture}
\newtheorem*{localconj*}{Local Conjecture}

\newtheorem*{nonzero*}{Conjecture on the non-vanishing of the normalized intertwining operators}
\newtheorem*{holo*}{Conjecture on the holomorphicity of the normalized intertwining operators}

\DeclareMathOperator{\ad}{ad}
\DeclareMathOperator{\Ad}{Ad}
\DeclareMathOperator{\supp}{supp}
\DeclareMathOperator{\vol}{vol}
\DeclareMathOperator{\grad}{grad}
\DeclareMathOperator{\tr}{tr}
\DeclareMathOperator{\Div}{div}
\DeclareMathOperator{\Hess}{Hess}
\DeclareMathOperator{\im}{im}
\DeclareMathOperator{\Aut}{Aut}
\DeclareMathOperator{\Hom}{Hom}
\DeclareMathOperator{\Gal}{Gal}
\DeclareMathOperator{\End}{End}
\DeclareMathOperator{\rank}{rank}
\newcommand{\trace}{\operatorname{trace}}

\newcommand{\U}{{\operatorname{U}}}
\renewcommand{\Res}{{\operatorname{Res}}}


\DeclareMathOperator{\abs}{abs}
\DeclareMathOperator{\exsupp}{ex.supp}

\let\oldbullet\bullet
\renewcommand{\bullet}{{\vcenter{\hbox{\tiny$\oldbullet$}}}}

\renewcommand{\1}{\mathbbm{1}}

\numberwithin{equation}{section}


\begin{document}

\title[induction restriction]{Parabolic induction and pullback of perverse sheaves on vogan varieties}



\author[Lo]{Chi-Heng Lo}
\address{Department of Mathematics,
National University of Singapore,
119076, Singapore}
\email{ch\_lo@nus.edu.sg}

\author[Ray]{Mishty Ray}
\address{Department of Mathematics\\
University of British Columbia\\
Vancouver, BC, V6T 1Z4,
CANADA}
\email{mishtyray@math.ubc.ca}


%    General info
\subjclass[2020]{Primary 11F70, 22E50; Secondary 11F85, 22E55}

\date{\today}

\keywords{Admissible Representations, Parabolic Induction, Local Arthur Packets, Local ABV Packets, Vogan Conjecture}

\thanks{}



\begin{abstract}
Let $M$ be a Levi subgroup of a quasisplit connected reductive group $G$ over a $p$-adic field. Vogan's geometric perspective on the local Langlands correspondence fixes a perfect pairing between the Grothendieck group of representations of $G$ and its pure inner forms with a fixed infinitesimal parameter and the Grothendieck group of equivariant perverse sheaves on the associated Vogan variety. Granting a list of desiderata for the enhanced local Langlands correspondence, we show that parabolic induction $\Ind_M^G$ is the transpose of geometric restriction $\xi^*$ with respect to this pairing if and only if the $p$-adic analogue of the Kazhdan--Lusztig hypothesis and the expected multiplicity formula for tempered induction both hold. The identity refines to one for each pure inner form of $G$. \end{abstract}

%Using this result, we formulate the notion of supercuspidal perverse sheaves and indicate applications to ABV-packets.}


\maketitle

\tableofcontents

\section{Introduction}\label{sec: introduction}


Let $F$ be a nonarchimedean local field of characteristic $0$ and let $G$ be a quasisplit connected reductive group over $F$.  The local Langlands correspondence predicts a finite-to-one parametrization of the smooth irreducible representations of $G(F)$ by $L$-parameters $\phi: W_F\times\SL_2(\BC)\to\Lgroup{G}$ where $\Lgroup{G}=\dualgroup{G}\rtimes W_F$ and $\dualgroup{G}$ is the $\CC$-points of the dual group of $G$.  An enhancement of the local Langlands correspondence refines the finite-to-one map into a bijection using irreducible representations of a component group attached to $\phi$. More precisely, on one side, we consider the set $\Pi^{\pure}(G)$ of equivalence classes of smooth irreducible representations of $G(F)$ along with those of its pure inner forms $G_\delta(F)$, as $\delta$ varies over elements of the Galois cohomology group $H^1(F,G)$. On the other hand, we have $\Phi^{\mfr{e}}(G)$, the set of equivalence classes of pairs $(\phi,\tau)$, where $\phi$ is a Langlands parameter and $\tau$ is an irreducible representation of the component group $A_\phi:=\pi_0(Z_{\dualgroup{G}}(\phi))$. This formulation is due to Vogan; see \cite{Vog93}, or for the archimedean case, \cite{ABV92}. We denote the enhanced local Langlands correspondence by
\[ \LLC^{\mfr{e}}: \Pi^{\pure}(G) \to \Phi^{\mfr{e}}(G).\]
The enhanced local Langlands correspondence is constructed for $\GL_n(F)$ in \cites{HT01, Hen00, scholze}, for classical groups in \cites{Art13, mokendoscopic2015, KMSW14, MR18, Ish24}, and for $G_2$ in \cites{GanSavin2023, aubert2023explicitlocallanglandscorrespondence}. There has been far more progress towards this conjecture than we are able to cite here; we move on for the sake of brevity.
A central aspect of this program is functoriality. An $L$-homomorphism $\Lgroup{H}\to\Lgroup{G}$ should give rise to a transfer of (packets of) representations from $H$ to $G$.  In \cite{ABV92}*{Proposition 1.37}, the authors formulate a Langlands functoriality map for archimedean groups in a geometric setting and utilize it when $H$ is an endoscopic group of $G$. This is extended to $p$-adic groups in \cite{Vog93} and reformulated in the language of Vogan varieties in \cite{CFMMX22}. We utilize the latter formulation in reviewing the geometric setting below. 

Given a Langlands parameter $\phi$, we attach an infinitesimal parameter by
\[ \lambda_{\phi}:= \phi\left(w,\begin{pmatrix}
    |w|^{1/2}&\\&|w|^{-1/2}
\end{pmatrix}\right).\]
The map $\phi \mapsto\lambda_{\phi}$ is a finite-to-one surjection modulo equivalence.  Let $\Phi_{\lambda}(G)$ be the fibre of this surjection at $\lambda$. We set $\Phi^{\mathfrak{e}}_\lambda(G)$ to be the subset of $\Phi^{\mfr{e}}(G)$ consisting of $(\phi,\tau)$ such that $\lambda_\phi=\lambda$ (or its $\dualgroup{G}$-conjugate). This enables us to start building the geometric spaces we need for Vogan's perspective.
Fixing an infinitesimal parameter $\lambda$ of $G$, we obtain the prehomogeneous vector space $V_\lambda$ called the {Vogan variety} with an action of the group $H_\lambda$ (see Section \ref{subsec: a moduli space for langlands parameters} for definitions).  Every $L$-parameter $\phi$ with infinitesimal parameter $\lambda$ determines an $H_\lambda$-orbit $C_\phi$, the representation $\tau$ determines an equivariant local system $\locsys{\tau}{C_\phi}$ on $C_\phi$, and $\mathcal{IC}(\locsys{\tau}{C_\phi})$ is the corresponding simple object of the category $\Per_{H_\lambda}(V_\lambda)$ of $H_\lambda$-equivariant perverse sheaves on $V_\lambda$. The enhanced local Langlands correspondence gives the bijections:
  \[\begin{array}{ccccc}
\Pi_\lambda^\pure(G) &{\longrightarrow} &\Phi^{\mathfrak{e}}_\lambda(G) & {\longrightarrow} &  \Per_{H_\lambda}(V_\lambda)^{\mathrm{simple}}_{\mathrm{/iso}} \\
\pi & {\longmapsto} & (\phi, \tau) & {\longmapsto} & \mathcal{IC}(\locsys{\tau}{C_\phi})=: \mathcal{P}(\pi).
\end{array}
\]
The above bijection defines a perfect pairing of Grothendieck groups
\[
    K\Pi^\pure_\lambda(G) \times K\Per_{H_\lambda}(V_\lambda) \to \mathbb{Z},
\]
\begin{equation}\label{equation: pairing in the introduction}
\langle \pi,\mathcal{P}\rangle =
\begin{cases}
    (-1)^{d(\pi)}e(\pi) & \text{if } \mathcal{P} = \mathcal{P}(\pi),\\
    0 & \text{otherwise},
\end{cases}
\end{equation}
where $d(\pi) \coloneqq \dim C_{\phi}$ and $e(\pi)$ is the Kottwitz sign of the pure inner form associated to $\pi$; see \eqref{equation: kottwitz sign}.

Returning to the question of functoriality, in this paper we only consider $H=M$, a Levi subgroup of $G$. Let $\xi: {}^L M \hookrightarrow {}^L G$ be the $L$-embedding. Fix an infinitesimal parameter $\lambda_M$ of $M$ and put $\lambda=\xi\circ\lambda_M$. We have the  $\Z$-linear map
\[\Ind_M^G:K\Pi^\pure_{\lambda_M}(M)\to K\Pi^\pure_{\lambda}(G),\qquad \pi \mapsto \Ind_P^G(\pi),\]
where $P$ is the standard parabolic subgroup of $G$ with Levi component $M$ and $\Ind_P^G$ is normalized parabolic induction (for a representation of a pure inner form of $M$ one induces to the corresponding pure inner form of $G$; see Section \ref{subsec: parabolic induction of reps linear map}). The map $\xi$ induces an inclusion $H_{\lambda_M} \hookrightarrow H_\lambda$ and an equivariant inclusion $V_{\lambda_M} \hookrightarrow V_\lambda$. We obtain the restriction functor $\xi^*:D_{H_\lambda}(V_\lambda) \to D_{H_{\lambda_M}}(V_{\lambda_M})$ between the equivariant bounded derived categories of constructible sheaves, and thus a $\Z$-linear map
\[\xi^*: K\Per_{H_\lambda}(V_\lambda) \to K\Per_{H_{\lambda_M}}(V_{\lambda_M})\]
between the respective Grothendieck groups. Then we define ${}^t\xi^* : K\Pi^\pure_{\lambda_M}(M)\to K\Pi^\pure_{\lambda}(G)$ by
\[ \langle {}^t\xi^*(\pi), \mathcal{P}\rangle=\langle \pi, \xi^*(\mathcal{P})\rangle\]
for any $\pi\in K\Pi_{\lambda_M}^\pure(M)$ and any $\mathcal{P}\in K\Per_{H_\lambda}(V_\lambda)$.
In this paper we establish the equality
 \begin{align}\label{eq ind res intro}
     \Ind_M^G={}^t\xi^*,
 \end{align}
 assuming specific desiderata from the local Langlands correspondence. Below, we expand on two of them, since they are of special importance to a converse statement. 

The first concerns the decomposition of tempered induction. If $(\phi_M, \tau_M)$ is an enhanced $L$-parameter of $M$ with $\phi_M$ tempered, and $\phi=\xi\circ\phi_M$, then one expects
\begin{align}\label{eq temp decomp intro}
\Ind_P^G\big(\pi(\phi_M,\tau_M)\big)=\bigoplus_{\tau\in \irrepdual{A}_\phi} m(\tau_M,\tau)\,\pi(\phi,\tau),
\end{align}
where $m(\tau_M,\tau)$ is the multiplicity of $\tau_M$ in the restriction of $\tau$ to $A_{\phi_M}$. This is Desideratum \ref{desideratum: induction of tempered representations}. 

The second is the $p$-adic analogue to the Kazhdan-Lusztig hypothesis. The Grothendieck group $K\Pi_\lambda^\pure(G)$ (resp. $K\Per_{H_\lambda}(V_\lambda)$) has two bases, one consists of irreducible representations $\pi(\phi,\tau)$ (resp. simple perverse sheaves $\mathcal{IC}(\locsys{\tau}{C_\phi})$), and the other consists of standard modules $\Delta(\pi)$ of $\pi$ (resp. standard sheaves $\stsh{\tau}{C_{\phi}}$). We let $m_{\lambda}$ and $c_{\lambda}$ be the change of basis matrices between the simple and standard objects in $K\Pi_\lambda^\pure(G)$ and $K\Per_{H_\lambda}(V_\lambda)$, respectively. The $p$-adic Kazhdan-Lusztig hypothesis states that 
\[ m_{\lambda}= {}^tc_{\lambda}.\]
Equivalently, the evaluation of the pairing \eqref{equation: pairing in the introduction} on standard objects is given by 
\begin{equation}\label{equation: pairing on standards in the introduction}
    \begin{array}{rcl}
        \langle \Delta(\pi), \stsh{\tau}{C_{\phi}} \rangle 
            &\mapsto& 
            \begin{cases}
                e(\pi) & \LLC^{\mathfrak{e}}_G(\pi) = (\phi,\tau), \\
                0 & \text{else},
            \end{cases}
    \end{array}
\end{equation}
where $e(\pi)$ is the Kottwitz sign for the pure inner form $\pi$ belongs to. This is Hypothesis \ref{hypothesis: KLH}.

%Let $m_\lambda$ be the matrix obtained by computing the multiplicity of an irreducible representation in a standard module, and let $c_{\lambda}$ be the matrix obtained by computing the Euler characteristic of the stalk at a chosen basepoint $x\in C$ of shifted simple perverse sheaves as $C$ varies over orbits each in $V_\lambda$. These matrices can be interpreted as change of basis matrices between the simple and standard objects in $K\Pi_\lambda^\pure(G)$ and $K\Per_{H_\lambda}(V_\lambda)$. This is the $p$-adic analogue of the Kazhdan--Lusztig hypothesis, going back to Zelevinsky for $\GL_n$. 

Conversely, the equality \eqref{eq ind res intro} implies the above two expectations. Here is our main theorem.

\begin{thm}[{Propositions \ref{proposition: pure parabolic induction and restriction}, \ref{theorem: ind res adjunction implies klh}}]\label{theorem: main theorem in the introduction}
    Let $G$ be a quasisplit connected reductive group over $F$. Assume the enhanced local Langlands correspondence for $G$ and for its standard Levi subgroups, together with the following desiderata:
    \begin{enumerate}
        \item compatibility with the Langlands classification, Desideratum \ref{desideratum: Langlands classification};
        \item compatibility with isomorphisms, Desideratum \ref{desideratum: isomorphism};
        \item compatibility of parabolic induction with infinitesimal parameters, Desideratum \ref{desideratum: Haines}.
    \end{enumerate}
    Then the following two statements are equivalent.
    \begin{enumerate}
        \item[(a)] The induction of tempered representations, Desideratum \ref{desideratum: induction of tempered representations}, and the $p$-adic analogue of the Kazhdan--Lusztig hypothesis, Hypothesis \ref{hypothesis: KLH}, both hold for $G$ and for each of its standard Levi subgroups.
        \item[(b)] For every standard Levi subgroup $L \leq M \leq G$ with $\xi_L^M: \dualgroup{L}\hookrightarrow \dualgroup{M}$, every infinitesimal parameter $\lambda_L$ of $L$ with $\lambda_M=\xi_{L}^M\circ\lambda_L$, all $\pi\in K\Pi_{\lambda_L}^\pure(L)$ and all $\mathcal{F} \in K\Per_{H_{\lambda_M}}(V_{\lambda_M})$,
        \[\langle \Ind_{L}^M(\pi),\mathcal{F}\rangle = \langle \pi,(\xi_{L}^M)^*(\mathcal{F})\rangle.\]
        In other words, $\Ind_L^M=\prescript{t}{}{(\xi_{L}^M)^*}$ with respect to the pairing \eqref{equation: pairing in the introduction}. 
    \end{enumerate}
\end{thm}



\subsection{Outline of the proof}
It suffices to treat the case of a Levi subgroup $M$ of $G$ with $\xi: \Lgroup{}{M}\hookrightarrow \Lgroup{G}$. For (a)$\implies$(b) in this case, we wish to show for all $\mathcal{F}\in K\Per_{H_\lambda}(V_\lambda)$
\[\langle \Ind_M^G(\pi), \mathcal{F}\rangle = \langle \pi, \xi^*(\mathcal{F})\rangle.\] Both sides of the identity 
%in Theorem \ref{theorem: main theorem in the introduction}(b) 
are linear, so it suffices to compare them on bases of standard objects in  $K\Pi_{\lambda_M}^\pure(M)$ and $K\Per_{H_\lambda}(V_\lambda)$. On right hand side, we use the basis of standard sheaves $\stsh{\tau}{C}$ of $K\Per_{H_\lambda}(V_\lambda)$, on which the restriction is computed in \cite{ABV92}*{Proposition 23.7(b)}; see \eqref{eqn: restriction formula on standards}:
\[
    \xi^*(\stsh{\tau}{C}) = \sum_{d\xi(D)\subseteq C,~\sigma \in \irrepdual{A}_D} m(\sigma,\tau) \cdot  \stsh{\sigma}{D}.
\]
Here, $D$ runs over the $H_{\lambda_M}$-orbits in $V_{\lambda_M}$ and $m(\sigma,\tau)$ is the multiplicity of $\sigma$ in $\tau \vert_{A_D}$. On the left hand side, we decompose $\Ind_M^G(\pi)$ using the basis of standard modules. From the decomposition \eqref{eq temp decomp intro}, we prove in Theorem \ref{theorem: induction formula} that parabolic induction decomposes standard modules in exactly the same way in the Grothendieck group:
\[\Ind_M^G(\Delta(\phi_M,\tau_M))=\sum_{\tau \in \irrepdual{A}_\phi} m(\tau_M,\tau)\, \Delta(\phi,\tau).\]
Then the $p$-adic analogue of the Kazhdan-Lusztig hypothesis says precisely that the pairing is diagonal on these two bases up to Kottwitz signs, which gives us the claim.

%Comparing the two formulas, and using that Kottwitz signs are compatible with Levi subgroups, gives (a)$\Rightarrow$(b) and its form-by-form refinement (Section \ref{subsection: pure parabolic induction and restriction} and Section \ref{subsection: parabolic induction form-by-form and lift}). 

For the converse, let $\Ind_M^G(\pi_M)$ be either {a tempered induction or a standard module}. In both cases, the representation $\pi_M$ has an $L$-parameter $\phi_M$ whose corresponding $H_{\lambda_M}$-orbit in $V_{\lambda_M}$ is open, and the question becomes the evaluation of 
\[\langle \Ind_M^G(\pi_M), \stsh{\tau}{C_{\phi}} \rangle.\]
where $\phi=\xi \circ \phi_M$  and $\tau_{C_\phi}$ is the extension by zero of the local system corresponding to $\tau$ on $C_\phi$. Assuming \eqref{eq ind res intro}, this is equal to 
\[  \langle \pi_M,  \xi^{\ast} \stsh{\tau}{C_{\phi}} \rangle. \]
 Then again using \cite{ABV92}*{Proposition 23.7(b)}, the desired conclusion follows.





\subsection{Related work and applications}
Theorem \ref{theorem: main theorem in the introduction} is a compatibility statement. It is the Levi case of the ``Langlands functoriality map'' of \cite{ABV92}*{Definition 1.36}, which Adams, Barbasch and Vogan use to describe the expected behaviour of their geometric packets in the archimedean case. The present result is the $p$-adic counterpart for Levi subgroups, where it has a clean interpretation in terms of parabolic induction.

For $G=\GL_n$, the identity of Theorem \ref{theorem: main theorem in the introduction}(b) was proved in \cite{CunninghamRayVoganGLn}. In this case the component groups $A_\phi=1$, $H^1(F,\GL_n)=1$, and parabolic induction from a tempered representation of a Levi subgroup is irreducible, so that Desideratum \ref{desideratum: induction of tempered representations} holds automatically. The $p$-adic Kazhdan--Lusztig hypothesis for $\GL_n$ (which was long regarded as a folklore theorem) is proved by Solleveld \cite{Sol25} for substantial classes of groups and representations including all classical groups; see also the exposition of Balodis in \cite{balodis2026proofpadickazhdanlusztighypothesis}. The identity is a key input in the proof of Vogan's conjecture for $p$-adic $\GL_n$ in \cite{CunninghamRayVoganGLn}: for an Arthur parameter $\psi$ of $\GL_n$, the Arthur packet and the ABV-packet coincide, $\Pi_\psi(G)=\Pi^\ABV_{\phi_\psi}(G)$, where $\phi_\psi$ is the Langlands parameter attached to $\psi$. This conjecture appears in \cite{ABV92} for archimedean groups and in \cite{Vog93} for $p$-adic groups, and it is formulated in the language of Vogan varieties and vanishing cycles in \cite{CFMMX22}*{Conjecture 1}.

Theorem \ref{theorem: main theorem in the introduction}(b) is related to geometric realizations of the Langlands-Shelstad and Kottwitz-Shelstad transfers (that is, maps of Grothendieck groups of perverse sheaves whose transposes give these transfers). These geometric realizations constitute a key step in proving Vogan's conjecture for groups for which endoscopic classification is known. This is taken up in \cite{CHLLRX26}. Specifically, see Theorem 15.3 of \textit{loc. cit.}; an important difference being that there the `transfer' and `lift' maps originate from the Grothendieck group of stable distributions on a quasisplit endoscopic group, whereas in this paper we work with the full Grothendieck group and do not restrict ourselves to the quasisplit case of Levi subgroup. 

Combined with the Fourier transform on $V_\lambda$ as in \cite{CFMMX22}*{Section 10.2.4}, our result is a step towards a geometric description of the Aubert-Zelevinsky involution for general connected reductive $p$-adic groups: the involution is compatible with parabolic induction, and the Fourier transform is compatible with the restriction $\xi^*$. The identity can also be used to obtain a neat characterisation of perverse sheaves corresponding to supercuspidal representations. Indeed, let $\mathcal{P}(\pi) \in \Per_{H_\lambda}(V_\lambda)_{/\text{iso}}^{\mathrm{simple}}$. Then $\pi$ is supercuspidal if and only if $\xi^*\mathcal{P}(\pi)=0$ for all  proper Levi subgroups $M <G$ for which $\lambda$ factors through the embedding $\xi:{}^L M \hookrightarrow {}^L G$ as $\lambda_M$, and $\xi^*:K\Per_{H_\lambda}(V_\lambda)\to K\Per_{H_{\lambda_M}}(V_{\lambda_M})$ is the corresponding equivariant pullback. This is treated in Section \ref{section: supercuspidal sheaves}.



\subsection{Acknowledgments} 
The authors would like to thank Clifton Cunningham for his constant support.
The second named author would like to thank Julia Gordon at University of British Columbia for providing a conducive research environment during which this paper was written.
\subsection{AI declaration}
We used LLMs to edit the document for typos, and as a search and summarization tool. All the ideation and mathematical work in this document is done by humans.
\iffalse
%{Let $F$ be a nonarchimedean local field of characteristic $0$ and let $G$ be a connected reductive quasisplit algebraic group defined over $F$. Throughout, $G$ is accompanied by all of its pure inner forms $G_\delta$, $\delta \in H^1(F,G)$, and every statement below about the whole family; the quasisplit form plays no privileged role beyond fixing the parametrization. The parabolic induction functor is central to study of admissible representations of $G$. If a local Langlands correspondence exists, two identities should be satisfied by this functor. If a local Langlands correspondence exists, then every invariant of a smooth irreducible representation of $G(F)$ ought to be legible in its enhanced $L$-parameter. This paper is concerned with one such invariant: the decomposition of a parabolically induced representation.}

The first concerns induction from a tempered representation of a Levi subgroup $M$ of $G$. Let $P$ be the standard parabolic subgroup of $G$ with Levi component $M$ and $\Ind_P^G$ denote normalised parabolic induction. Let $\xi: \Lgroup{M} \xhookrightarrow{}\Lgroup{G}$ denote the inclusion of $L$-groups. Here the expectation is purely group-theoretic: if $(\phi_M, \tau_M)$ an enhanced $L$-parameter of $M$ where $\tau_M$ is an irreducible representation of the component group $A_{\phi_M}=\pi_0(Z_{\dualgroup{M}}(\phi_M))$, and $\phi=\xi\circ\phi_M$ is the corresponding parameter of $G$, then
\[{\Ind_P^G\big(\pi(\phi_M,\tau_M)\big)=\bigoplus_{\tau\in \irrepdual{A}_\phi} m(\tau_M,\tau)\,\pi(\phi,\tau),}\]
{Here $m(\tau_M,\tau)$ is the multiplicity of $\tau_M$ in the restriction of $\tau$ to the component group $A_{\phi_M}$. This is Desideratum \ref{desideratum: induction of tempered representations} below. It is expected to follow from the endoscopic character identities, see Remark \ref{remark: tempered induction from ECI}.}

The second concerns the decomposition of a standard module of a representation. Here the expected answer is not group-theoretic but geometric: the multiplicity of an irreducible representation in a standard module should be computed by the multiplicity of the corresponding simple perverse sheaf in the corresponding standard sheaf, on a variety of Langlands parameters attached to the infinitesimal parameter. This is the $p$-adic analogue of the Kazhdan-Lusztig hypothesis, going back to Zelevinsky for $\GL_n$ and to Vogan \cite{Vog93} in general; it appears as Hypothesis \ref{hypothesis: KLH} below.

{In this paper, assuming the enhanced local Langlands classification and a short list of desiderata, i.e, \ref{desideratum: Langlands classification} and \ref{desideratum: isomorphism}, the two aforementioned conjectures together are equivalent to a single geometric description of parabolic induction: parabolic induction of representations is the transpose (with respect to a pairing \eqref{equation: pairing in the introduction} of Grothendieck groups) of the restriction with of constructible sheaves along the inclusion induced from $\xi$ of the corresponding varieties of Langlands parameters. This is Theorem \ref{theorem: main theorem in the introduction}. 
%One direction packages the two conjectures into one clean statement; the other direction says that nothing is lost in the packaging, since the geometric statement recovers both.
}

The framework in which these statements are made involve the geometric perspective on the local Langlands correspondence for $p$-adic groups due to Vogan in \cite{Vog93}. This perspective involves relating representations to perverse sheaves on a moduli space of Langlands parameters as we now explain. Fix an infinitesimal parameter $\lambda:W_F \to \prescript{L}{}{G}$ of $G$, defined in Section \ref{subsec: a moduli space for langlands parameters}. Vogan's enhancement of the local Langlands correspondence gives us the identification
\[
\begin{array}{ccccc}
\Pi_\lambda^\pure(G) &{\longrightarrow} &\Phi^{\mathfrak{e}}_\lambda(G) & {\longrightarrow} &  \Per_{H_\lambda}(V_\lambda)^{\mathrm{simple}}_{\mathrm{/iso}} \\
(\pi,\delta) & {\longmapsto} & (\phi, \tau) & {\longmapsto} & \mathcal{IC}(\locsys{\tau}{C_\phi}) 
\end{array}
\]
where the arrows are all bijections and the objects are explained as follows.
\begin{itemize}
    \item The set $ {\Pi^\pure(G)}=\bigsqcup_{\delta\in H^1(F,G)}\Pi(G_\delta)$ consists of equivalence classes of smooth irreducible representations of $G$ along with its pure inner forms; the latter are indexed by $\delta\in H^1(F,G)$.
    \item The set $\Phi^{\mathfrak{e}}_\lambda(G)$ consists of equivalence classes of pairs $(\phi,\tau)$, where $\phi:W_F\times \SL_2(\CC)\to \Lgroup{G}$ is a Langlands parameter whose infinitesimal parameter is $\lambda$, and $\tau$ is an irreducible representation of the component group $A_\phi=\pi_0(Z_{\dualgroup{G}}(\phi))$ of $\phi$. 
    \item $ {\Pi^\pure_\lambda(G)}$ is the preimage of $ {\Phi^{\mathfrak{e}}_\lambda(G)}$ in $ {\Pi^\pure(G)}$ under the enhanced local Langlands correspondence.
    \item The space $V_\lambda$ is an affine variety with the action of the group $H_\lambda$ as explained in Section \ref{subsec: a moduli space for langlands parameters}. Every $L$-parameter determines an $H_\lambda$-orbit $C_\phi$. The element $\tau$ corresponds to a local system $\locsys{\tau}{C_\phi}$ on $C_\phi$ and its intersection cohomology complex $\mathcal{IC}(\locsys{\tau}{C_\phi})$ gives an element of  the set $\Per_{H_\lambda}(V_\lambda)^{\mathrm{simple}}_{\mathrm{/iso}}$ of equivalence classes of $H_\lambda$-equivariant simple perverse sheaves on $V_\lambda$. 
\end{itemize}
The enhanced Langlands correspondence comes equipped with a list of desiderata; see for example  {\cite{Bor79}, \cite{Vog93}, \cite{Kal16} and \cite{AMS18}}. We assume a minimal list of desiderata for our context, as listed below in the statement of Theorem \ref{theorem: main theorem in the introduction}.
%\todo[inline]{Chi-Heng: Do you know a good place where these desiderata are listed?\\ \textbf{[disposition:} Borel (Corvallis), Vogan, Kaletha, Aubert--Moussaoui--Solleveld now cited; two bib keys still to add.\textbf{]}}
%\revnote{\rev{the keys \texttt{Bor79} (Borel, \emph{Automorphic $L$-functions}, Corvallis~II) and \texttt{AMS18} (Aubert--Moussaoui--Solleveld, \emph{Generalizations of the Springer correspondence and cuspidal Langlands parameters}) still have to be added to \texttt{references.bib}, and precise section pointers inserted here.}} 

We drop the $\delta$ from the notation and write an element of $\Pi_\lambda^\pure(G)$ as $\pi$. The map $\pi \mapsto \mathcal{P}(\pi):=\mathcal{IC}(\locsys{\tau}{C_\phi})$ fixes a perfect pairing between Grothendieck groups
\[
    K\Pi^\pure_\lambda(G) \times K\Per_{H_\lambda}(V_\lambda) \to \mathbb{Z}
\]
defined by
\begin{equation}\label{equation: pairing in the introduction}
\langle [\pi],[\mathcal{P}]\rangle =
\begin{cases}
    (-1)^{d(\pi)}e(\pi) & [\mathcal{P}] = [\mathcal{P}(\pi)],\\
    0 & \text{otherwise},
\end{cases}
\end{equation}
where $d(\pi) \coloneqq \dim C_{\phi}$, and $e(\pi):=e(\pi(\phi,\tau))$ is the Kottwitz sign; see  {\eqref{equation: kottwitz sign}}. 

Let $M$ be a Levi subgroup of $G$. Let $\lambda$ be an infinitesimal parameter of $G$ that factors through the parameter $\lambda_M$ of $M$, that is, $\lambda=\xi\circ\lambda_M$, where $\xi:\Lgroup{M}\hookrightarrow \Lgroup{G}$. We write down the linear map 
\[\Ind_M^G:K\Pi^\pure_{\lambda_M}(M)\to K\Pi^\pure_{\lambda}(G)\] 
via $\pi \mapsto \Ind_P^G(\pi)$, where $P=MN$ is a standard parabolic subgroup of $G$ and $\Ind_P^G$ is the normalised parabolic induction functor. The map $\xi$ also gives us the inclusion $H_{\lambda_M} \hookrightarrow H_\lambda$ and the equivariant inclusion $V_{\lambda_M} \hookrightarrow V_\lambda$ with respect to the $H_{\lambda_M}$-action (resp. $H_\lambda$-action) on $V_{\lambda_M}$ (resp. $V_\lambda$). This gives us the linear map
\[\xi^*: K\Per_{H_\lambda}(V_\lambda) \to K\Per_{H_{\lambda_M}}(V_{\lambda_M})\]
via $\mathcal{F} \mapsto \xi^*(\mathcal{F})$. The map $\xi^*:D_{H_\lambda}(V_\lambda) \to D_{H_{\lambda_M}}(V_{\lambda_M})$ is the equivariant restriction of sheaves where $D_H(V)$ is the 
$H$-equivariant bounded derived category of constructible sheaves on $V$.
Our main result establishes the natural identity relating $\Ind_M^G$ and $\xi^*$:
\begin{thm}\label{theorem: main theorem in the introduction}
    Let $G$ be a  {quasisplit} connected reductive algebraic group over a $p$-adic field $F$ {, let $M$ be a Levi subgroup of $G$, and let $\lambda=\xi\circ\lambda_M$ be as above}.  {Assume} the enhanced local Langlands correspondence  {for $G$ and for its Levi subgroups, together with the two desiderata}
    \begin{enumerate}
        \item compatibility with Langlands classification {, Desideratum} \ref{desideratum: Langlands classification} {;}
        \item compatibility with isomorphism {, Desideratum} \ref{desideratum: isomorphism} {.}
        \end{enumerate}
     {Then the following two statements are equivalent.}
    \begin{enumerate}
        \item[ {(a)}]  {The induction of tempered representations, Desideratum \ref{desideratum: induction of tempered representations}, and the $p$-adic analogue of the Kazhdan--Lusztig hypothesis, Hypothesis \ref{hypothesis: KLH}, both hold.}
        \item[ {(b)}]  {For all $\pi\in K\Pi_{\lambda_M}^\pure(M)$ and all $\mathcal{F} \in K\Per_{H_\lambda}(V_\lambda)$,}
        \[\langle \Ind_M^G(\pi),\mathcal{F}\rangle = \langle \pi,\xi^*(\mathcal{F})\rangle;\]
         {in other words, $\Ind_M^G=\prescript{t}{}{\xi^*}$ with respect to the pairing \eqref{equation: pairing in the introduction}.}
    \end{enumerate}
    Moreover,  {when these hold,} one can refine  {the identity in (b)}: let $\gamma \in H^1(F,M)$ and suppose $\delta=H_\xi(\gamma)$ where $H_\xi: H^1(F,M) \to H^1(F,G)$ is defined in \eqref{equation: definition of h xi}. Then, for any $\pi \in K\Pi_{\lambda_M}(M_\gamma)$ and $\mathcal{F}\in K\Per_{H_\lambda}(V_\lambda)^\delta$ (see \eqref{equation: kper subspace delta}),
    \[\langle \Ind_{M_\gamma}^{G_\delta}(\pi),\mathcal{F} \rangle = \langle \pi, 
    \xi^*(\mathcal{F})\rangle.\]
     {In particular the theorem is not a statement about the quasisplit form alone: this second identity is the statement of \textup{(b)} for the Levi subgroup $M_\gamma$ of the pure inner form $G_\delta$, and Proposition \ref{lemma: lift is induction} rewrites it as}
    \[ {\Ind_{M_\gamma}^{G_\delta}=\Lift_{M_\gamma}^{G_\delta}=\pr_\delta\circ \prescript{t}{}{\xi^*}\circ \iota_\gamma.}\]
\end{thm}
\subsection{Outline of proof} {The implication (a)$\Rightarrow$(b) is Proposition \ref{proposition: pure parabolic induction and restriction}, and the converse (b)$\Rightarrow$(a) is Proposition \ref{theorem: ind res adjunction implies klh} together with Proposition \ref{corollary: converse tempered}; the refinement is Proposition \ref{lemma: lift is induction}.} The proof  {of (a)$\Rightarrow$(b)} involves understanding the maps $\Ind_M^G$ and $\xi^*$ with respect to the basis of standard objects {:} standard representations in $K\Pi_{\lambda_M}^\pure(M)$ and constructible standard sheaves in $K\Per_{H_\lambda}(V_\lambda)$. As before, we have an infinitesimal parameter $\lambda: W_F\to\Lgroup{G}$  factoring through the homomorphism $\lambda_M: W_F\to\Lgroup{M}$. 
%\todo[inline]{check exactly which desiderate is used in which step\\ \textbf{[disposition:} done --- the three desiderata are now named explicitly here and in Theorem \ref{theorem: induction formula}.\textbf{]}}
Using Desiderata  {\ref{desideratum: Langlands classification}, \ref{desideratum: induction of tempered representations} and \ref{desideratum: isomorphism}}, one computes  {in Theorem \ref{theorem: induction formula} the effect of $\Ind_M^G$ on the standard module $\Delta(\phi_M,\tau_M)$ of $\pi(\phi_M,\tau_M)\in \Pi_{\lambda_M}^\pure(M)$}: 
\[ {\Ind_M^G(\Delta(\phi_M,\tau_M))=\sum_{\tau \in \irrepdual{A}_\phi} m(\tau_M,\tau)\, \Delta(\phi,\tau).}\]
Here  {$m(\tau_M,\tau)$ is the multiplicity of $\tau_M$ in $\tau\vert_{A_{\phi_M}}$}; see \eqref{equation: pullback of tau}. On the other hand, suppose $C$ is an $H_\lambda$-orbit in $V_\lambda$. From \eqref{eqn: restriction formula on standards},  {the} equivariant restriction of a  {standard} $H_{\lambda}$-equivariant constructible sheaf  {$\stsh{\tau}{C}$}, corresponding to $\tau \in \irrepdual{A}_C$, is given by 
\[
    \xi^*( {\stsh{\tau}{C}}) = \sum_{ {d\xi(D)\subseteq C},~\sigma \in \irrepdual{A}_D} m(\sigma,\tau) \cdot  {\stsh{\sigma}{D}},
\]
where $D$ varies over $H_{\lambda_M}$-orbits in $V_{\lambda_M}$, $\sigma \in \irrepdual{A}_D$, and $m(\sigma,\tau)$ is the multiplicity of $\sigma$ in $\tau \vert_{A_D}$.  {Comparing the two formulas against the pairing yields}
\[\langle \Ind_M^G(\pi), \mathcal{F}\rangle = \langle \pi,  {\xi^*(\mathcal{F})}\rangle\] for any $\pi \in K\Pi_{\lambda_M}^\pure(M)$ and $\mathcal{F}\in K\Per_{H_\lambda}(V_\lambda)$; in other words, $\Ind_M^G=\prescript{t}{}{\xi^*}$.

Now, let $M_\gamma$ be a pure inner form of $M$ where $\gamma\in H^1(F,M)$ and  {$\delta:=H_\xi(\gamma)\in H^1(F,G)$}. We define the map
 {$\Ind_{M_\gamma}^{G_\delta}:K\Pi_{\lambda_M}(M_\gamma)\to K\Pi_\lambda(G_\delta)$} via $\pi \mapsto \Ind_{P}^{G_\delta}(\pi)$ where $P$ is a parabolic subgroup of $G_\delta$ containing $M_\gamma$ as a Levi subgroup {; such a $P$ exists precisely because $\delta$ is in the image of $H_\xi$, so that $M_\gamma$ is indeed a Levi subgroup of $G_\delta$ (see \eqref{equation: definition of h xi})}. %\todo[inline]{Clarify this ?? Or clarify if we need to clarify this :/ !!!\\ \textbf{[disposition:} clarified here and at the end of Section \ref{subsection: setting}: $\delta$ determines the fibre $H_\xi^{-1}(\delta)$, not a single $\gamma$.\textbf{]}}
Then one may define the map $\Lift_{M_\gamma}^{G_\delta}$ as the composition 
\[K\Pi_{\lambda_M}(M_\gamma) \xhookrightarrow{\iota_\gamma} K\Pi_{\lambda_M}^\pure(M) \xrightarrow[]{\prescript{t}{}{\xi^*}} K\Pi_{\lambda}^\pure(G)\xrightarrow[]{\pr_{\delta}} K\Pi_{\lambda}(G_\delta).\]
We see that for any $\pi \in K\Pi_{\lambda_M}(M_\gamma)$
\[\Ind_{M_\gamma}^{G_\delta}(\pi)=\Ind_M^G(\iota_\gamma(\pi))=\Lift_{M_\gamma}^{G_\delta}(\pi).\]
 {For a fixed pair of forms one cannot simply write $\Ind_{M_\gamma}^{G_\delta}$ as $\prescript{t}{}{\xi^*}$, because $\xi^*$ need not ``preserve forms'': if $\stsh{\tau}{C}$ is attached to the class $\delta \in H^1(F,G)$, then every $\stsh{\sigma}{D}$ occurring in $\xi^*(\stsh{\tau}{C})$ is attached to {some} class $\gamma \in H^1(F,M)$ with $H_\xi(\gamma)=\delta$ (Lemma \ref{lemma: relating galois colomology class of irreps of component groups}), and $H_\xi$ need not be injective. This is exactly why the projection $\pr_\delta$ appears in the definition \eqref{equation: lift} of $\Lift_{M_\gamma}^{G_\delta}$. The Kottwitz signs are nonetheless matched, $e(\gamma)=e(\delta)$, by  \eqref{eqn Kottwitz sign Levi}, and it is this compatibility that makes the refined identity of Theorem \ref{theorem: main theorem in the introduction} hold form by form.}

%\old{[\textsf{\textbf{superseded, for comparison.}} The previous version of this passage read: ``For a fixed form, one cannot write $\Ind_{M_\gamma}^{G_\delta}$ as $\prescript{t}{}{\xi^*}$ since $\xi^*$ does not \emph{preserve forms}. Example ?? illustrates this: if $\tau_C$ corresponds to the quasisplit form, that is, $e(\kappa_G(\tau))=1$, then $\xi^*(\tau_C)$ may contain $\sigma_D$ where $e(\kappa_M(\sigma))=-1$.'' The last assertion contradicts Lemma \ref{lemma: galois cohomology map}, which gives $e(\gamma)=e(H_\xi(\gamma))$ for every $\gamma \in H^1(F,M)$: the Kottwitz sign is always preserved. What genuinely fails is that $H_\xi$ need not be injective, so the \emph{class} $\gamma$, and hence the form, is not determined by $\delta$. The $\SO_7$ example of Section \ref{section: example} does not exhibit even that, since there $H_\xi$ is bijective.]}
\todo[inline]{ Replace all \texttt{cite[A]{B}} by \texttt{cite{B}*{A}}}
{\subsection{Related work and applications}
Theorem \ref{theorem: main theorem in the introduction} is a compatibility statement.
It is the Levi case of the `Langlands functoriality map' that appears in \cite[Definition 1.36]{ABV92}; Adams, Barbasch, and Vogan use this framework to describe the expected behaviour of their geometric packets in the archimedean case. The present result supplies the $p$-adic counterpart along an $L$-embedding $\xi: \Lgroup{M}\hookrightarrow \Lgroup{G}$ specialized to Levi subgroups with a neat relationship with parabolic induction.

For $G=\GL_n$ the identity of Theorem \ref{theorem: main theorem in the introduction}(b) is a theorem due to Cunningham and Ray \cite{CunninghamRayVoganGLn}. In this special case, component groups $A_\phi$ are trivial and relatedly, $H^1(F,\GL_n)=1$. Parabolic induction from a tempered representation of a Levi subgroup of $\GL_n$ is irreducible so that Desideratum \ref{desideratum: induction of tempered representations} holds automatically. The $p$-adic Kazhdan--Lusztig hypothesis for $\GL_n$ was considered a folklore theorem but is also proven  in Solleveld \cite{Sol25} where he verifies it for substantial classes of groups and representations, including all classical groups. This identity proves Vogan's conjecture for $p$-adic $\GL_n$, that is, given an Arthur parameter $\psi$ of $\GL_n$, the Arthur packet and the ABV-packet coincide: $\Pi_\psi(G)=\Pi^\ABV_{\phi_\psi}(G)$ where $\phi_\psi$ is the corresponding Langlands parameter of $\psi$. This conjecture appears for archimedean groups in \cite{ABV92} and for the $p$-adic case in \cite{Vog93}. It is reformulated in the language of Vogan varieties in \cite[Conjecture 1]{CFMMX22}. 

A generalization of the identity \ref{theorem: main theorem in the introduction}(b) involves geometric realizations of (that is, maps of Grothendieck groups of perverse sheaves whose transpose gives) Langlands-Shelstad transfer and Kottwitz-Shelstad transfer and is a key part of the strategy to prove Vogan's conjecture for a larger class for groups for which Arthur's endoscopic classification is known. This is taken up in upcoming work by the authors of this paper in collaboration with Cunningham, Hazeltine, Liu and Xu. 

Finally, together with the Fourier transform on $V_\lambda$ as in \cite[Section 10.2.4]{CFMMX22}, it is a first step towards a geometric description of the Aubert-Zelevinsky involution for a general connected reductive $p$-adic group, since the involution is compatible with parabolic induction and the Fourier transform is compatible with the restriction $\xi^*$. \fi


\section{Background from representation theory}
Throughout, $F$ is a nonarchimedean local field of characteristic $0$ with residue field of size $q=p^f$, for a fixed prime $p$ and an integer $f$.
\subsection{Groups and representations}\label{Groups and representations}
Let $G$ be a quasisplit connected reductive algebraic group. Fix a Borel subgroup $B \subset G$ and a maximal torus $T \subseteq B$. 
Let \[X^*(T)=\operatorname{Hom}(T,\mathbb{G}_m),\qquad X_*(T)=\operatorname{Hom}(\mathbb{G}_m,T)\] denote the character and cocharacter lattices of $T$ respectively. 
Let $\Sigma \subset X^*(T)$
be the root system for $(G,T)$ and $\dualgroup{\Sigma} \subset X_*(T)$
the corresponding coroot system. The choice of $B$ determines a set $\Sigma^+$ of positive roots
and a corresponding set $\Delta$ of simple roots. For each simple root $\alpha\in\Delta$, let
$
X_\alpha:\mathbb{G}_a\longrightarrow U_\alpha
$
be a fixed isomorphism onto the corresponding root subgroup. The data
$
(G,B,T,\{X_\alpha\}_{\alpha\in\Delta})
$
is called a \emph{pinning} of $G$. The pinning determines a canonical splitting of the based root datum
$
(X^*(T),\Sigma,\Delta,X_*(T),\dualgroup{\Sigma},\dualgroup{\Delta}),
$
and identifies the group of pinned automorphisms of $G$ with the automorphism group of the based root datum. Throughout this paper, all constructions are understood with respect to this fixed pinning. Moreover, since we fix a Borel subgroup of $G$, we also get a collection of standard parabolic subgroups of $G$. Associated to the pinned group $(G,B,T,\{X_\alpha\})$ is the pinned Langlands dual group
$(\dualgroup{G},\dualgroup{B},\dualgroup{T},\{X^\vee_{\alpha^\vee}\}),
$ whose based root datum is dual to that of $G$. In this paper, we consider the complex points $\dualgroup{G}(\CC)$ of the dual group of $G$. {We often shorten the notation of the complex points of the dual group to simply $\dualgroup{G}$.}


Next we fix notation for the Weil group. Fix an algebraic closure $\overline{F}$ of $F$. The Weil group is a dense subgroup of the absolute Galois group $\Gal(\overline{F}/F)$ that satisfies the short exact sequence
\[1 \rightarrow I_F\rightarrow W_F \rightarrow \mathbb{Z} \rightarrow 1\]
where $I_F$ is the inertia subgroup and
$\Z$ is generated by the geometric Frobenius map $\Fr$ of $F$. We let $|\cdot|$ denote the absolute value on $W_F$ such that $|\Fr|=q$.

 The action of the absolute Galois group
$
\Gamma_F=\operatorname{Gal}(\overline{F}/F)
$
on the based root datum of $G$ induces an action on $\dualgroup{G}$ by pinned automorphisms. The $L$-group of $G$ is the semidirect product
\[
{}^LG = \dualgroup{G}\rtimes W_F.
\]
For any standard Levi subgroup $M$ of $G$, there is an embedding between their $L$-groups
\begin{align}\label{eq Levi L embedding}
    \xi: {}^L M \hookrightarrow {}^L G,
\end{align} well-defined up to conjugation by $\dualgroup{G}$. See \cite{Bor79}*{Section 15.1} or \cite{Tai25}*{Section 4.3} for details.  



Finally, we fix notation for the two sets of representations that occur throughout. For a finite group $A$, let $\irrepdual{A}$ denote the set of equivalence classes of complex irreducible representations of $A$. 
We let $\Pi(G)$ denote the set of equivalence classes of complex smooth irreducible representations of $G(F)$. We let $K\Pi(G)$ denote the Grothendieck group of smooth representations of finite length of $G$.

\subsection{Parabolic induction}\rev{\label{subsec: parabolic induction}}
In this section, $G$ is a connected reductive algebraic group, not necessarily quasisplit.  For a parabolic subgroup $P$ of $G$ with a Levi subgroup $M$, we let $\Ind_P^G$ denote the normalized parabolic induction from finite length representations of $M(F)$ to representations of $G(F)$. More precisely, let $(\sigma,W)$ be an admissible representation of $M$ which is inflated to a representation of $P$. Then $\Ind_P^G(\sigma)=(\pi,V)$, where $V$ is the $\CC$-vector space of smooth functions $f:G(F)\to W$ such that 
\[f(pg)=\delta_P(p)^{1/2}\sigma(p)f(g) ~\text{ for }~ p\in P(F) ~\text{ and }~g\in G(F),\]
where $\delta_P$ is the modular quasicharacter of $P$. The $G$-action on $V$ is given by $\pi(g)f(x)=f(xg)$ for all $g,x\in G(F)$. 

Suppose that $M$ is also a Levi factor of another parabolic subgroup $P'$ of $G$. It is known that $\Ind_P^G(\sigma)$ and $\Ind_{P'}^G(\sigma)$ have the same image in $K\Pi(G)$ (\cite{BZ80}*{Theorem 2.9, Remark 2.10}). Therefore, we let 
\[ \Ind_M^G:K\Pi(M) \to K\Pi(G) \]
denote the linear transformation  $\sigma \mapsto \Ind_P^G(\sigma)$ for any choice of $P$.
\subsection{The enhanced local Langlands correspondence}\label{subsection: the enhanced local Langlands correspondence}
A Langlands parameter for $G$ is a continuous homomorphism $\phi: W_F \times \SL_2(\BC) \to \Lgroup{G}$ satisfying conditions specified in \cite{CFMMX22}*{Section 3.4}. An $L$-parameter is a $\dualgroup{G}$-conjugacy class of a Langlands parameter. We write $\phi$ for both the Langlands parameter and the $L$-parameter, unless it is necessary to emphasize the distinction. We let $\Phi(G)$ denote the set of $L$-parameters of $G$. 

In its classical form, the local Langlands correspondence is a conjectural many-to-one surjection between the set $\Pi(G)$ of smooth irreducible representations of $G(F)$ and $L$-parameters of $G$:
 \[ \LLC: \Pi(G) \to \Phi(G).\]
The fibres of $\LLC$ are the $L$-packets. In our setting, the conjectural surjection is refined to a bijection using the component group for each $\phi$,
\begin{equation}\label{eqn: a phi}
    A_\phi = Z_{\dualgroup{G}}(\phi)/Z_{\dualgroup{G}}(\phi)^0,
\end{equation}
where $Z_{\dualgroup{G}}(\phi)$ denotes the centralizer of the image of $\phi$ in the dual group $\dualgroup{G}$, and $Z_{\dualgroup{G}}(\phi)^0$ is its identity component. 
Let $\Phi^\mathfrak{e}(G)$ denote the set of $\dualgroup{G}$-conjugacy classes of pairs $(\phi, \tau)$, where $\phi$ is a Langlands parameter for $G$ and $\tau \in \irrepdual{A}_\phi$. Such a pair is called an \emph{enhanced $L$-parameter}. If we restrict our attention to the quasisplit group $G$ alone, the fiber of $\LLC$ over $\phi$ is parametrized by the quotient $A_{\phi}/Z(\dualgroup{G})^{\Gamma_F}$. However, to incorporate Vogan's geometric perspective (see Section~\ref{sec: geometric LLC}), it is more natural to work with $A_{\phi}$ itself. This viewpoint naturally leads us to consider the pure inner forms of $G$ simultaneously, which are parametrized by the Galois cohomology group $H^1(F,G)$. Let
\[\Pi^\pure(G)=\bigsqcup_{\delta\in H^1(F,G)}\Pi(G_\delta)\]
denote the set of equivalence classes of smooth irreducible representations of the quasisplit group $G$, which corresponds to the trivial element in $ H^1(F,G)$,
along with its other pure inner forms $G_{\delta}$. The enhanced local Langlands correspondence is a canonical bijection 
\begin{equation}\label{eqn: enhanced llc}
    \begin{array}{cc}
             \LLC^\mathfrak{e}: \Pi^\pure(G) &\to \Phi^\mathfrak{e}(G),  \\
         \pi&\mapsto (\phi, \tau). 
    \end{array}
\end{equation}
We often write $\pi(\phi,\tau)$ to denote the representation $\pi$ which satisfies $\LLC^\mathfrak{e}(\pi)=(\phi,\tau)$. 
We let $\LLC$ denote the composition of $\LLC^\mathfrak{e}$ with the projection $(\phi, \tau) \mapsto \phi$. Then we define $L$-packets
\[ \Pi_{\phi}^{\pure} (G):= \LLC^{-1}(\phi),\ \ \Pi_{\phi} (G_{\delta}):= \Pi_{\phi}^{\pure} (G) \cap \Pi(G_{\delta}).\]
The bijection \eqref{eqn: enhanced llc} is expected to depend on a choice of Whittaker datum $\mathfrak{w}$ of $G$, which we have fixed. We write $\LLC^{\mfr{e}}_{\mfr{w}}$ for $\LLC^{\mfr{e}}$ and $\pi_{\mfr{w}}(\phi,\tau)$ for $\pi(\phi,\tau)$ to emphasize the dependence on $\mfr{w}$ if necessary. 

The class $\delta$ attached to an enhancement is read off by the Kottwitz isomorphism, which we now recall. For a Langlands parameter $\phi$, the inclusion
\[Z(\dualgroup{G})^{\Gamma_F} \hookrightarrow Z_{\dualgroup{G}}(\phi)\] { induces the well-defined map of their component groups}
\begin{equation}\label{equation: inclusion of componenet grp of center in a phi}
    \iota_\phi: \pi_0(Z(\dualgroup{G})^{\Gamma_F}) \to A_\phi.
\end{equation}
Moreover, the image of $\iota_\phi$ is contained in the center of $A_\phi$. For any $\tau \in \irrepdual{A}_\phi$, the pullback $\tau \circ \iota_\phi$ is an isotypic direct sum of the pullback of the central character of $\tau$ to $\pi_0(Z(\dualgroup{G})^{\Gamma_F})$. We let $\chi_{\tau} \in \Hom(\pi_0(Z(\dualgroup{G})^{\Gamma_F}), \BC^\times)$ be the corresponding character. Recall the Kottwitz isomorphism for $G$:
\begin{equation}\label{eqn: kottwitz map}
    \kappa_G: \Hom(\pi_0(Z(\dualgroup{G})^{\Gamma_F}), \BC^\times) \xrightarrow[]{\sim} H^1(F,G).
\end{equation}
For $\tau \in \irrepdual{A}_\phi$, define 
\begin{equation}\label{eqn: kappa tau}
    \kappa_G(\tau):=\kappa_G(\chi_{\tau}).
\end{equation}

Often, we drop the subscript and use $\kappa(\tau)$ instead of $\kappa_G(\tau)$ when the group is clear. Finally, we build the following requirement into our definition of $\LLC^{\mathfrak{e}}$: for $(\phi,\tau)\in \Phi^{\mathfrak{e}}(G)$ and $\delta \in H^1(F,G)$,
\begin{equation}\label{equation: kappa determines the form}
    \kappa(\tau)=\delta \quad \Longleftrightarrow \quad \pi(\phi,\tau)\in \Pi(G_\delta).
\end{equation} We recall the definition of the Kottwitz sign.
Write $r(H)$ for the $F$-rank of the derived group of $H$. For $\delta \in H^1(F,G)$, the \emph{Kottwitz sign} of the pure inner form $G_\delta$ is
\begin{equation}\label{equation: kottwitz sign}
    e(\delta):= e(G_\delta)=(-1)^{r(G)-r(G_\delta)}.
\end{equation}
In this way, to each $\delta \in H^1(F,G)$, we assign the Kottwitz sign $\pm 1$. This may be interpreted as a function $e:H^1(F,G)\to \{\pm 1\}$. We can attach a Kottwitz sign to $\tau \in \Irrep(A_\phi)$ using \eqref{eqn: kappa tau}:
\begin{equation}\label{equation: kottwitz sign of tau}
    e(\tau):=e(\kappa_G(\tau)).
\end{equation}

With the requirement \eqref{equation: kappa determines the form} the enhanced local Langlands correspondence \eqref{eqn: enhanced llc} restricts to a bijection for each $\delta \in H^1(F,G)$: 
\begin{equation}\label{equation: llc delta}
    \LLC^{\mathfrak{e}}_{\delta}: \Pi(G_{\delta}) \to \{ (\phi, \tau) \in \Phi^\mathfrak{e}(G) \ \vert\  \kappa(\tau)=\delta \}.
\end{equation}



\subsection{Langlands classification}\label{subsec: Langlands classification}
We recall the Langlands classification, for representations and for $L$-parameters, together with its refinement to enhanced $L$-parameters.

Let $\delta\in H^1(F,G)$. We let $\Phi(G_{\delta})$ denote the subset of $\Phi(G)$ that is $G_{\delta}$-relevant. Namely, $\phi \in \Phi(G_{\delta})$ if and only if the minimal Levi subgroup of ${}^LG$ containing the image of $\phi$ is conjugate to the $L$-group of a standard Levi subgroup of $G_{\delta}$.

The Langlands classification for both $\Pi(G_{\delta})$ (\cite{Kon03}*{Theorem 3.5}) and $\Phi(G_{\delta})$ (\cite{SZ18}) provides canonical bijections
\begin{align*}
    \pi \leftrightarrow (P, \pi_\temp, \nu), \qquad \phi \leftrightarrow (P, \phi_\temp, \nu),
\end{align*}
where:
\begin{itemize}
    \item $\pi \in \Pi(G_{\delta})$ and $\phi \in \Phi(G_{\delta})$, 
    \item $P$ is a standard parabolic subgroup of $G_{\delta}$ with standard Levi component $M$,
    \item $\nu$ is a positive real unramified character of $M$ lying in the positive Weyl chamber determined by $P$, and
    \item $\pi_\temp \in \Pi(M)$ and $\phi_\temp \in \Phi(M)$ are both tempered (i.e., $\phi_\temp|_{W_F}$ has a bounded image).
\end{itemize}


%\todo[inline]{Check consistency for \texttt{emph} definition.}

Under this correspondence, $\pi$ is realized as the unique irreducible quotient of the \emph{standard module} (or \emph{standard representation})
\[
\Delta(\pi):=\Ind_P^{G_{\delta}}(\pi_\temp \otimes \nu),
\]
and the triple $(P,\pi_\temp,\nu)$ is uniquely determined by $\pi$.

For $L$-parameters, the positive real unramified character $\nu$ determines an unramified cocycle $c_{\nu}:W_F \to Z(\dualgroup{M})$. Define the twist $\phi_\temp \otimes \nu$  by 
\[(\phi_\temp \otimes \nu)(w,x):= c_{\nu}(w) \phi_{\temp}(w,x); \]
see \cite{SZ18}*{Definition 4.5}.
The relation between $\phi$ and the triple $(P, \phi_\temp, \nu)$ is given by
\[
\phi = ({}^L M \hookrightarrow {}^L G) \circ (\phi_\temp \otimes \nu).
\]
Since $\phi$ is $G_{\delta}$-relevant, the Levi subgroup $M$ is a Levi subgroup of $G_{\delta}$.



We need the following ingredients from the proof of the Langlands classification for $\Phi(G_{\delta})$.
\begin{prop}[{\cite{SZ18}*{Lemma 5.2, Propositions 5.5, 6.1, equation (6) in \S 7.2}}]\label{prop: SZ18}
   Let $\phi$ be a Langlands parameter of $G$. Suppose that the corresponding $L$-parameter $[\phi]$ corresponds to $(P, [\phi_\temp], \nu)$ under the Langlands classification. Then the following holds.
   \begin{enumerate}
    \item The hyperbolic part $z_{\phi}$ of $\phi(\Fr)$ is independent of the choice of $\Fr$. Moreover, the character $\nu$ is determined by the $\dualgroup{G}$-conjugacy class of $z_{\phi}$.
    \item We have ${}^LM $ conjugate to $ Z_{{}^L G}(z_{\phi})$ under $\dualgroup{G}$.
    \item The homomorphism of component groups $A_{\phi_\temp}^{\phi}(\xi): A_{\phi_\temp} \to A_{\phi}$ induced from the $L$-embedding $\xi: \prescript{L}{}{M} \hookrightarrow \prescript{L}{}{G}$ is an isomorphism.
   \end{enumerate}
\end{prop}

By Part (3) in the above proposition, the Langlands classification for $\Phi(G)$ can be further refined to the Langlands classification for $\Phi^{\mathfrak{e}}(G)$. More precisely, we define
\[ (\phi, \tau) \leftrightarrow (P, (\phi_\temp, \tau_\temp), \nu),\]
if $\phi \leftrightarrow (P, \phi_\temp, \nu)$ and $\tau_\temp = \tau \circ A_{\phi_\temp}^{\phi}(\xi)$.


It is also useful to define the set
\begin{equation}\label{eqn: a phi delta}
    \irrepdual{A}_{\phi,\delta}:=\{\tau \in \irrepdual{A}_\phi\ | \  \kappa(\tau)=\delta\},
\end{equation}
%One writes \todo[inline]{this is not lemma, but an obvious consequence of the desideratum?}
which allows us to write
\[\irrepdual{A}_\phi=\bigsqcup_{\delta\in H^1(F,G)}\irrepdual{A}_{\phi, \delta}.\]

\section{Background from geometry}\label{sec: geometric LLC}
The main goal of this section is to outline the geometric objects that inform Vogan's geometric perspective on the local Langlands correspondence from \cite{Vog93} which of course depends on the existence of $\LLC^{\mathfrak{e}}$. We interpret this geometric perspective using a pairing on Grothendieck groups.

\subsection{A moduli space for Langlands parameters}\label{subsec: a moduli space for langlands parameters}
An \emph{infinitesimal parameter} for $G$ is a continuous homomorphism $\lambda: W_F \to \Lgroup{G}$ that is a section of $\Lgroup{G} \to W_F$ and whose image consists of semisimple elements in $\Lgroup{G}$; see \cite{CFMMX22}*{Section 4.1}. The map $\phi \mapsto \lambda_\phi$ is surjective, and finite-to-one modulo equivalence.  We now fix an infinitesimal parameter $\lambda$ of $G$. The Vogan variety for $\lambda$ is 
\begin{equation}\label{equation: vogan variety}
    V_\lambda=\{x \in \mathrm{Lie}(\dualgroup{G})~ \vert ~ \Ad(\lambda(w))x=|w|x, ~\forall w\in W_F \}.
\end{equation}
It carries the $\Ad$-action of the centralizer group
\begin{equation}\label{equation: h lambda}
    H_\lambda=Z_{\dualgroup{G}}(\lambda).
\end{equation}
The complex affine variety $V_\lambda$ is a prehomogeneous vector space stratified under the $H_\lambda$-action into finitely many orbits. 

\subsection{Simple Perverse Sheaves}\label{subsec: simple perverse sheaves}
For an $H_\lambda$-orbit $C$ in $V_\lambda$, choose a base point $x \in C$. The equivariant fundamental group of $C$ is the finite group $A_C=\pi_0(Z_{H_\lambda}(x))$. Let $\Per_{H_\lambda}(V_\lambda)$ denote the category of $H_\lambda$-equivariant perverse sheaves on $V_\lambda$. It is the heart of the perverse $t$-structure on the $H_\lambda$-equivariant bounded derived category $D_{H_\lambda}(V_\lambda)$ of constructible sheaves on $V_\lambda$. Equivalence classes of simple equivariant local systems on $C$ are parametrized by the set $\irrepdual{A}_C$ of equivalence classes of irreducible representations of the finite group $A_C$. Let $\locsys{\tau}{C}$ denote the local system on $C$ corresponding to $\tau \in \irrepdual{A}_C$. The set $\Per_{H_\lambda}(V_\lambda)_{\text{/iso}}^{\text{simple}}$ of simple objects up to equivalence in $\Per_{H_\lambda}(V_\lambda)$ is given by
\begin{equation}\label{eq: simple perverse sheaves classified}
    \Per_{H_\lambda}(V_\lambda)_{\text{/iso}}^{\text{simple}}=\{\mathcal{IC}(\locsys{\tau}{C})~\vert ~ H_\lambda \text{-orbit }C \subseteq V_\lambda, \tau \in \irrepdual{A}_C \}.
\end{equation}
Here $\mathcal{IC}(\locsys{\tau}{C})$ denotes 
the corresponding intersection cohomology complex: given the inclusions $C\xhookrightarrow{j}\ol{C}\xrightarrow[]{i}V_\lambda$, 
\begin{equation}\label{equation: ic definition}
    \mathcal{IC}(\locsys{\tau}{C}):=i_*j_{!*}\locsys{\tau}{C}[\dim C]
\end{equation}
where $j_{!*}$ is the intermediate extension functor. See \cite{CFMMX22}*{Chapter 4}, or for more general exposition on perverse sheaves and derived categories for use in representation theory, see \cite{Ach21}*{Sections  3.1, 3.3, 6.4}. 

\subsection{Standard sheaves}\label{subsec: standards sheaves} The $H_\lambda$-equivariant bounded derived category $D_{H_\lambda}(V_\lambda)$ has another heart, the category $\operatorname{Sh}_{H_\lambda}(V_\lambda)$ of $H_\lambda$-equivariant constructible sheaves on $V_\lambda$; see \cite{Ach21}*{Section 6.4, Exercise 6.4.3}. We explain the objects of this category. Let $C$ be an $H_\lambda$-orbit in $V_\lambda$. For every local system $\locsys{\tau}{C}$ supported on $C$, we obtain a constructible sheaf on $V_\lambda$ via the extension by zero functor (see \cite{Ach21}*{Section 1.3}):
\begin{equation}
    \stsh{\tau}{C} :=k_!\locsys{\tau}{C} \text{ where } k:C \hookrightarrow V_\lambda.
\end{equation} 
We refer to the sheaves $\stsh{\tau}{C}$ as standard sheaves. The simple objects up to equivalence are
\begin{equation}\label{eq: standard sheaves classified}
    \operatorname{Sh}_{H_\lambda}(V_\lambda)_{\text{/iso}}^{\text{simple}}=\{\stsh{\tau}{C}~\vert ~ H_\lambda \text{-orbit }C \subseteq V_\lambda, \tau \in \irrepdual{A}_C \}.
\end{equation}


\subsection{Vogan--Langlands correspondence}\label{Vogan--Langlands correspondence}
We recall Vogan's perspective on the local Langlands correspondence for the quasisplit connected reductive algebraic group $G$ and its pure inner forms, established in \cite{Vog93}. As explained in \cite{CFMMX22}*{Proposition 4.2}, every point on the variety $V_\lambda$ corresponds to a Langlands parameter with infinitesimal parameter $\lambda$. They are related by 
\[ \phi \mapsto x_{\phi}=d(\phi|_{\SL_2})\begin{pmatrix}
    0 & 1\\0& 0
\end{pmatrix} \in V_{\lambda}. \]
Every $H_\lambda$-orbit in $V_\lambda$ corresponds to a conjugacy class of such a Langlands parameter: there is a bijection $\phi \mapsto C_{\phi}$ between $L$-parameters of $G$ that share the infinitesimal parameter $\lambda$, and the $H_{\lambda}$-orbits $C_{\phi}$ in $V_{\lambda}$. 
From
\cite{CFMMX22}*{Lemma 4.4}, for each Langlands parameter $\phi$,
we have that \begin{equation}\label{eqn: equivalence of component group and eq fund group}
    A_{C_{\phi}}\simeq A_{\phi}.
\end{equation}
Recall that the definition of $A_{C_{\phi}}$ depends on a basepoint in $C_{\phi}$, which we choose to be $x_{\phi}$ for the above isomorphism.

Let $\Pi_{\lambda}^\pure(G)$ denote the subset of representations in $\Pi^\pure(G)$ whose Langlands parameters share the infinitesimal parameter $\lambda$. Let $\pi(\phi,\tau)$ be the representation satisfying $\LLC^{\mathfrak{e}}(\pi(\phi,\tau))=(\phi, \tau)$, where $\tau \in \irrepdual{A}_\phi$; see \eqref{eqn: enhanced llc}. Vogan's geometric perspective on the local Langlands correspondence can be summarized as the bijection
\begin{equation}\label{eqn: vogan langlands correspondence}
    \begin{array}{rcl}
    \Pi_{\lambda}^\pure(G) 
    &\leftrightarrow& \Per_{H_\lambda}(V_\lambda)_{\text{/iso}}^{\text{simple}} \\
    \pi(\phi, \tau) 
    &\mapsto&
    \mathcal{IC}(\locsys{\tau}{C_{\phi}})=:\mathcal{P}(\pi).
\end{array}
\end{equation}
\subsection{Grothendieck groups}\label{subsec: grothendieck groups}
Let $KD_{H_\lambda}(V_\lambda), K\Per_{H_\lambda}(V_\lambda), $ and $ K\operatorname{Sh}_{H_\lambda}(V_\lambda)$ denote the Grothendieck groups of the categories $D_{H_\lambda}(V_\lambda)$, $\Per_{H_\lambda}(V_\lambda)$, and $\operatorname{Sh}_{H_\lambda}(V_\lambda)$, respectively. One can interpret $K\Per_{H_\lambda}(V_\lambda)$ (resp.\ $K\operatorname{Sh}_{H_\lambda}(V_\lambda)$) as the $\Z$-span of the elements of the set \eqref{eq: simple perverse sheaves classified} (resp. \eqref{eq: standard sheaves classified}).

As in \cite{ABV92}*{Equation 7.10(a)}, we have the identification
\[K\Per_{H_\lambda}(V_\lambda)=K\operatorname{Sh}_{H_\lambda}(V_\lambda)=KD_{H_\lambda}(V_\lambda).\]
In particular, $K\Per_{H_\lambda}(V_\lambda)$ can be interpreted as a $\mathbb{Z}$-module (or, after a change of scalars, a $\BC$-module) with two bases: one consisting of simple perverse sheaves as in \eqref{eq: simple perverse sheaves classified} and the other consisting of standard constructible sheaves as in \eqref{eq: standard sheaves classified},
\begin{equation}\label{equation: gr grp perverse sheaves}
\begin{aligned}
    K\Per_{H_\lambda}(V_\lambda)&=\mathrm{Span}_{\mathbb{Z}}\{\mathcal{IC}(\locsys{\tau}{C}): C \text{ an }H_\lambda\text{-orbit in }V_\lambda, \tau \in \irrepdual{A}_C \}\\
    &=\mathrm{Span}_{\mathbb{Z}}\{\stsh{\tau}{C}: C \text{ an }H_\lambda\text{-orbit in }V_\lambda, \tau \in \irrepdual{A}_C \} {.}
\end{aligned}
\end{equation}
The orbits can be labeled using Langlands parameters by invoking the Vogan-Langlands correspondence. For example, recalling $\irrepdual{A}_{\phi,\delta}$ from \eqref{eqn: a phi delta}, we define
\begin{equation}\label{equation: kper subspace delta}
    K\Per_{H_\lambda}(V_\lambda)^\delta := \operatorname{Span}_{\mathbb{Z}} \{\mathcal{IC}(\locsys{\tau}{C_\phi}): \lambda_\phi=\lambda \text{ and } \tau \in  \irrepdual{A}_{\phi,\delta}\} {.}
\end{equation}
{We will see the representation theory counterpart of the two bases in \eqref{equation: gr grp perverse sheaves} in Section \ref{subsec: Kazhdan--Lusztig hypothesis}.}



We record a case of special interest that relates the extension by zero of a local system with its intersection cohomology complex. Suppose $C$ is an open $H_\lambda$-orbit in $V_\lambda$, that is, $\ol{C}=V_\lambda$. Keeping \eqref{equation: ic definition} in mind, we write $j: C \hookrightarrow{} V_\lambda$ and $\mathcal{IC}(\locsys{\tau}{C})=j_{!*}\locsys{\tau}{C}[\dim C]$. Note $\locsys{\tau}{C}[\dim C] \in \Per(C)$. We compare $\stsh{\tau}{C}=j_!\locsys{\tau}{C}$ and $\mathcal{IC}(\locsys{\tau}{C})$ in $K\Per_{H_\lambda}(V_\lambda)=KD_{H_\lambda}(V_\lambda)$: 
\[
j_!\locsys{\tau}{C}[\dim C]
=
\mathcal{IC}(\locsys{\tau}{C})
+
\sum_{\substack{D\subsetneq \ol{C}\\ \gamma \in \irrepdual{A}_D}}
m(D,\gamma)
\mathcal{IC}(\locsys{\gamma}{D}),
\]
where $m(D,\gamma)\in\mathbb{Z}$. One can deduce this relation using \cite{BBD82}*{p. 54-55}. Shifting by an integer $d$ is the same as multiplication by $(-1)^d$ in the Grothendieck group. We thus obtain that the standard sheaf $\stsh{\tau}{C}$, up to the important sign $(-1)^{\dim C}$, is a sum of the perverse sheaf $\mathcal{IC}(\locsys{\tau}{C})$ and sheaves supported on the boundary $V_\lambda\setminus C$:
\begin{equation}\label{equation: standard sheaf as a sum of perverse sheaf and boundary terms}
    (-1)^{\dim C} \stsh{\tau}{C}=\mathcal{IC}(\locsys{\tau}{C}) + \sum_{\substack{D\subsetneq \ol{C}\\ \gamma \in \irrepdual{A}_D}}
m(D,\gamma)
\mathcal{IC}(\locsys{\gamma}{D})
\end{equation}
for some integers $m(D,\gamma)$.
\subsection{Pairing of Grothendieck groups on the basis of simple objects}\label{subsec: Pairing simple}
For every infinitesimal parameter $\lambda :W_F \to \Lgroup{G}$, the Vogan-Langlands correspondence, as in \eqref{eqn: vogan langlands correspondence}, fixes a perfect pairing between Grothendieck groups
\[
    K\Pi^\pure_\lambda(G) \times K\Per_{H_\lambda}(V_\lambda) \to \mathbb{Z}
\]
defined by
\begin{equation}\label{eqn:pairing}
\langle \pi,\mathcal{P}\rangle =
\begin{cases}
    (-1)^{d(\pi)}e(\pi) & \mathcal{P} = \mathcal{P}(\pi),\\
    0 & \text{otherwise},
\end{cases}
\end{equation}
where $\LLC^{\mathfrak{e}}(\pi)=(\phi,\tau)$, $d(\pi) \coloneqq \dim C_{\phi}$, and $e(\pi)$ is the Kottwitz sign attached to $\pi$, that is $e(\pi):=e(\delta)$ (see \eqref{equation: kottwitz sign}) where $\pi \in \Pi(G_\delta)$. This pairing is equivalent to the one in \cite{Vog93}*{(8.11)'(a)}; subsequently it appears in \cite{CFMMX22}*{p. 77} using the language of Vogan varieties. 
\section{Desiderata}
We gather the list of desiderata of the enhanced local Langlands correspondence needed in this work.

\subsection{Compatibility with Langlands classification}\label{subsec: Compatibility with Langlands classification}
Recall from Section \ref{subsec: Langlands classification} the Langlands classification for representations and $L$-parameters. We consider not only the local Langlands correspondence for pure inner forms of $G$, but also the local Langlands correspondence for pure inner forms of all standard Levi subgroups of $G$. These correspondences are compatible via Langlands classification in the following sense.
\begin{desideratum}\label{desideratum: Langlands classification}
Suppose that $\LLC^{\mathfrak{e}}(\pi) = (\phi, \tau)$. Let 
\[\pi \leftrightarrow (P_{\pi}, \pi_{\mathrm{temp}}, \nu_{\pi}),\ \ \ (\phi, \tau) \leftrightarrow (P_{\phi}, (\phi_{\mathrm{temp}}, \tau_{\mathrm{temp}}), \nu_{\phi})\] 
under the Langlands classification. Then we have $P_{\pi} = P_{\phi}$ and $\nu_{\pi} = \nu_{\phi}$, and furthermore,
\[
    \LLC^{\mathfrak{e}}(\pi_{\mathrm{temp}}) = (\phi_{\mathrm{temp}}, \tau_{\mathrm{temp}}).
\]
\end{desideratum}

\begin{remark}\label{rmk LC tempered}
Desideratum \ref{desideratum: Langlands classification} implies that if $\pi$ is essentially tempered, then its $L$-parameter $\phi= \phi_{\mathrm{temp}} \otimes \nu_{\phi}$ is a twist of a tempered $L$-parameter $\phi_{ \mathrm{temp}}$. In particular, since $\phi_{\mathrm{temp}}$ corresponds to the open orbit in the associated Vogan variety by \cite{cunningham2024whittakernormalizationpadicabvpackets}*{Proposition 3.2}, so does $\phi$.
\end{remark}


\subsection{Induction of tempered representations}\label{subsec: Induction of tempered representations}
The second desideratum involves the decomposition of the parabolic induction $\Ind_P^G (\pi(\phi_M,\tau_M))$ when $\pi(\phi_M,\tau_M)$ is a tempered representation of the Levi subgroup $M$ of $G$. It is known that this induction is semisimple and consists of tempered representations of $G$. We describe the subquotients of this induction in terms of the enhanced $L$-parameters of $G$.

First, we introduce some notation. Note that we do not assume $\phi$ is tempered in this paragraph.
Let $\xi: \prescript{L}{}{M} \hookrightarrow \prescript{L}{}{G}$ be the $L$-embedding. We let $\phi := \xi \circ \phi_M$. The embedding $\xi$ induces a homomorphism between the component groups of $\phi_M$ and $\phi$:
\[ A_{\phi_M}^{\phi}(\xi): A_{\phi_M} \to A_{\phi}.\]
For any $\tau \in \irrepdual{A}_{\phi}$, we decompose the pullback of $\tau$ along $A_{\phi_M}^{\phi}(\xi)$ into a direct sum of irreducible representations of $A_{\phi_M}$ as
\begin{align}\label{equation: pullback of tau}
\tau \circ A_{\phi_M}^{\phi}(\xi) = \bigoplus_{\tau_M \in \irrepdual{A}_{\phi_M}} m(\tau_M, \tau)\tau_M, 
\end{align}
where $m(\tau_M, \tau) \in \Z_{\geq 0}$ represents the multiplicity of $\tau_M$ in this decomposition. We now use this to describe the decomposition of the tempered induction. 

\begin{desideratum}\label{desideratum: induction of tempered representations}
    Let $\phi_M$ be a tempered $L$-parameter of $M$ and $\tau_M \in \irrepdual{A}_{\phi_M}$. Then the parabolic induction $\Ind_P^G (\pi(\phi_M,\tau_M))$ decomposes as
    \begin{align}\label{eq decomp of tempered induction}
        \Ind_{P}^{G} (\pi(\phi_M,\tau_M))=\bigoplus_{\tau \in \irrepdual{A}_{\phi}} m(\tau_M, \tau) \pi(\phi, \tau).
    \end{align}
\end{desideratum}

\begin{remark} {\label{remark: tempered induction from ECI}}
    A stronger desideratum is to require that the tempered $L$-packet satisfies endoscopic character identities $($see \cite{Kal16}*{Conjecture E}$)$. Indeed, by the compatibility of endoscopic transfer and parabolic induction $($\cite{Hir04}*{Lemma 6.1}$)$, the desired equation \eqref{eq decomp of tempered induction} can be obtained by a straightforward computation using the Fourier inversion formula. A detailed argument will be provided in \cite{KLT}.
\end{remark}

In Theorem \ref{theorem: induction formula} we extend the identity \eqref{eq decomp of tempered induction} to the standard modules of nontempered representations in the Grothendieck group.

\subsection{Compatibility with isomorphism}\label{subsection: compatibility with isomorphism}
 We fix an isomorphism of quasisplit algebraic groups $f: G_1 \to G_2$ defined over $F$ throughout this subsection. It induces a bijection $ H^1(F,G_1) \to H^1(F,G_2)$ and we identify these two sets. In particular, for any $\delta \in H^1(F,G_1)$, the isomorphism $f$ also induces an isomorphism
$f_{\delta}: G_{1,\delta} \simeq G_{2, \delta}$. Moreover, we also have a bijection
\begin{align*}
    \Pi(G_{2,\delta}) & \to \Pi(G_{1,\delta}),\\
    \pi_2 & \mapsto \pi_2 \circ f_{\delta}.
\end{align*}



 Let ${}^L f: {}^L G_1 \to {}^L G_2$ be the corresponding $L$-isomorphism. It induces a bijection $\Phi(G_1) \to \Phi(G_2)$ given by $\phi_1 \mapsto {}^L f \circ \phi_1$, and also an isomorphism $ A_{\phi_1} \simeq A_{{}^L f \circ \phi_1}$, which we also denote by ${}^L f$. Thus, we obtain a map between enhanced $L$-parameters
 \begin{align*}
     \Phi^{\mfr{e}}(G_{1}) &\to  \Phi^{\mfr{e}}(G_2),\\
     (\phi_1,\tau_1) & \mapsto ({}^L f \circ \phi_1, \tau_1 \circ ({}^Lf)^{-1} )=: {}^Lf (\phi_1,\tau_1).
 \end{align*}


We need a desideratum that the enhanced local Langlands correspondence on $G_1$ and $G_2$ is compatible with the above bijections. 
\begin{desideratum}\label{desideratum: isomorphism}
    
   Fix a Whittaker datum $\mfr{w}_2=(B_2, \chi_2)$ of $G_2$ and let $\mfr{w}_1:= (f^{-1}(B_2), \chi_2 \circ f)$ be the corresponding Whittaker datum of $G_1$. The enhanced local Langlands correspondences for $G_1$ and $G_2$ satisfy the following: for all $(\phi_1, \tau_1) \in \Phi^{\mfr{e}}(G_1)$ with $\kappa_{G_1}(\tau_1)= \delta$, we have 
    \[ \pi_{\mfr{w}_1}(\phi_1,\tau_1)= \pi_{\mfr{w}_2}({}^Lf (\phi_1,\tau_1))   \circ f_{\delta}. \]
\end{desideratum}



We now describe the setting in which this desideratum will be applied in the
later section. Write $A_0$ for the maximal split subtorus of the maximal torus
$T$ of $G$ fixed in Section \ref{Groups and representations}. Let $M$ be a
standard Levi subgroup of $G$ and let $g \in N_G(A_0)(F)$ be an element such
that $M' := g M g^{-1}$ is again a standard Levi subgroup. Conjugation by $g$
then defines an isomorphism $f \colon M \to M'$, and the induced
$L$-isomorphism ${}^L f \colon {}^L M \to {}^L M'$ is given by conjugation by
some $\widehat{g} \in G^{\vee}$. For $\delta \in H^1(F,M)$ and
$\pi \in \Pi(M_{\delta})$ we write $\pi^g := \pi \circ f_{\delta}^{-1}$, where
$f_{\delta} \colon M_{\delta} \to M'_{\delta}$ is the isomorphism of pure inner
forms induced by $f$. For $(\phi, \tau) \in \Phi^{\mfr{e}}(M)$ we write
$({}^{\widehat{g}}\phi, {}^{\widehat{g}}\tau) := {}^L f (\phi,\tau)$.

Fix a Whittaker datum $\mfr{w} = (B,\chi)$ for $G$, and let $\mfr{w}_{M}$ and
$\mfr{w}_{M'}$ be the Whittaker data for $M$ and $M'$ obtained by restriction. Namely, write $N_B$ the unipotent radical of the Borel subgroup $B$, we have
\[ \mfr{w}_M = \bigl( B \cap M , \chi|_{N_B \cap M} \bigr), \]
and similarly for $M'$. Note that as long as $B$ is contained in a parabolic subgroup $P$ with Levi component $M$, the $M(F)$-conjugacy class of this restriction only depends on the $G(F)$-conjugacy class of $\mfr{w}$.

Now we check that $\mfr{w}_M$ and $\mfr{w}_{M'}$ are
related as in Desideratum \ref{desideratum: isomorphism}. Namely, the
pullback
\[ f^{*}\mfr{w}_{M'} := \bigl( f^{-1}(B \cap M'),\;
   \chi|_{N_B \cap M'} \circ f \bigr) \]
is equivalent to $\mfr{w}_M$. Indeed, $f^{*}\mfr{w}_{M'}$ is the restriction to
$M$ of the Whittaker datum $\mfr{w}' := (g^{-1}Bg,\; \chi \circ \Ad(g))$ of
$G$. Since $\mfr{w}'$ and  $\mfr{w}$ are conjugate by
$g \in G(F)$, their restriction is conjugate by $M(F).$
Hence, $\mfr{w}'_M = f^{*}\mfr{w}_{M'}$ is equivalent to
$\mfr{w}_M$. The desideratum therefore yields
\[ \LLC^{\mfr{e}}_{\mfr{w}_{M'}}
   \bigl( \pi_{\mfr{w}_M}(\phi,\tau)^g \bigr)
   = ({}^{\widehat{g}} \phi, {}^{\widehat{g}} \tau) \]
for every $(\phi, \tau) \in \Phi^{\mfr{e}}(M)$.

Finally, suppose we start instead from $\delta \in H^1(F,G)$, a standard Levi
subgroup $M_{\delta}$ of $G_{\delta}$, and an element $g_{\delta} \in
G_{\delta}(F)$ such that $M_{\delta}' := g_{\delta} M_{\delta} g_{\delta}^{-1}$
is again standard. Then one can choose $g \in N_G(A_0)(F)$ conjugating the
quasisplit Levi subgroup $M$ to $M'$ in such a way that the induced map
$f_{\delta} \colon M_{\delta} \to M_{\delta}'$ agrees with conjugation by
$g_{\delta}$. Then the discussion above applies.


\begin{remark}
    The description of $L$-parameter in Desideratum \ref{desideratum: isomorphism} is a special case of \cite{Bor79}*{Condition 10.3(5)}. The description of the enhanced $L$-parameters can be deduced from \cite{Kal23}*{Section 2.3.4}.
\end{remark}

%\todo[inline]{MR: I think we discussed moving KLH in the geometric background section; we can do that if you wish\\ \textbf{[disposition:} left here, but Section \ref{subsec: grothendieck groups} now sets up both transition matrices, so the move is no longer needed.\textbf{]}}
\subsection{\texorpdfstring{Haines conjecture and the $p$-adic Kazhdan--Lusztig hypothesis}{}}\label{subsec: Kazhdan--Lusztig hypothesis}
In this subsection, we discuss the decomposition of the standard modules in terms of local Langlands correspondence.

First, we state Haines' conjecture on compatibility of induction and infinitesimal parameter \cite{Hai14}*{Conjecture 5.2.2}.
\begin{desideratum}\label{desideratum: Haines}
Let $\lambda_M$ be an infinitesimal parameter of $M$, a Levi subgroup of $G$. Let $\lambda:= ({}^L M \hookrightarrow {}^L G) \circ \lambda_M$ be the corresponding infinitesimal parameter of $G$. Suppose that $\pi_M  \in \Pi^{\pure}_{\lambda_M}(M)$ and $\pi$ is an irreducible subquotient of $\Ind_M^G\pi_M$. Then $ \pi \in \Pi^{\pure}_{\lambda}(G)$.
\end{desideratum}




An important special case of the above desideratum is that if $\pi \in \Pi_{\lambda}^{\pure}(G)$, then $\Delta(\pi) \in K\Pi_{\lambda}^{\pure}(G)$. Since the change of basis matrix from standard modules to irreducible representations is unitriangular (see \cite{BW00}*{Lemma XI.2.13}), the set 
\[ \{\Delta(\pi)\ | \ \pi \in \Pi^{\pure}_{\lambda}(G) \}\]
forms a $\Z$-basis. Therefore, assuming Desideratum \ref{desideratum: Haines}, the Grothendieck group $K\Pi_\lambda^\pure(G)$ has two bases {:} one consisting of equivalence classes of irreducible representations $\pi \in \Pi_\lambda^\pure(G)$ and the other consisting of equivalence classes of their standard modules $\Delta(\pi)$:
\begin{equation}\label{equation: gr grp of reps}
    K\Pi_\lambda^\pure(G)=\mathrm{Span}_{\mathbb{Z}}~\Pi_\lambda^\pure(G)=\mathrm{Span}_{\mathbb{Z}}\{\Delta(\pi)\ | \  \pi \in \Pi_\lambda^\pure(G)\}.
\end{equation}
Similarly, for $\delta\in H^1(F,G)$,
\[K\Pi_\lambda(G_\delta)=\mathrm{Span}_{\mathbb{Z}}~\Pi_\lambda(G_\delta)=\mathrm{Span}_{\mathbb{Z}}\{\Delta(\pi)\ | \  \pi \in \Pi_\lambda(G_{\delta})\}.\]


The Kazhdan--Lusztig hypothesis relates the change of basis matrix from irreducible representations to standard modules in $K\Pi_\lambda^\pure(G)$ to the change of basis matrix from simple perverse sheaves to standard constructible sheaves in $K\Per_{H_\lambda}(V_\lambda)$. See Section \ref{subsec: grothendieck groups} for aforementioned bases of Grothendieck groups. The statement of the hypothesis agrees with \cite{Sol25}*{Theorem E} and with \cite{Vog93}*{Conjecture 8.11} after dropping signs in the latter. This sign correction 
is illustrated via a simple example in \cite{balodis2026proofpadickazhdanlusztighypothesis}*{Example 3.5}. We state the hypothesis below in terms of the evaluation of the pairing \eqref{eqn:pairing} on the bases of standard objects in $K\Pi_\lambda^\pure(G)$ and $K\Per_{H_\lambda}(V_\lambda)$. The relation between this formulation and the change of basis matrices is explained in \cite{CunninghamRayVoganGLn}*{Section 1.4}.

 

\begin{hypothesis}[Kazhdan--Lusztig Hypothesis for $G$]\label{hypothesis: KLH}
Desideratum \ref{desideratum: Haines} holds. Moreover, the pairing defined in \eqref{eqn:pairing} is equivalent to the pairing defined by
\begin{equation}\label{equation: pairing on standards}
    \begin{array}{rcl}
        K\Pi_\lambda^\pure(G) \times K\Per_{H_\lambda}(V_\lambda)
            &\to& \mathbb{Z}\\
        \langle \Delta(\pi), \stsh{\tau}{C_{\phi}} \rangle 
            &\mapsto& 
            \begin{cases}
                e(\pi) & \LLC^{\mathfrak{e}}_G(\pi) = (\phi,\tau), \\
                0 & \text{else},
            \end{cases}
    \end{array}
\end{equation}
where $\stsh{\tau}{C_\phi}$ denotes the standard sheaf corresponding to $\tau$ on the $H_\lambda$-orbit $C_\phi$, and $e(\pi)$ is the Kottwitz sign of $\pi$ as in \eqref{eqn:pairing}. 
\end{hypothesis}

\section{Parabolic induction and restriction as maps between Grothendieck groups}
We recall and prove identities for parabolic induction of representations and equivariant restriction of sheaves at the level of Grothendieck groups.

\subsection{Setting}\label{subsection: setting}
Let $G$ be a connected reductive quasisplit algebraic group over a nonarchimedean local field $F$. Let $M$ be a standard Levi subgroup and $P$ be the standard parabolic subgroup with Levi component $M$. Recall from \eqref{eq Levi L embedding} that we have an $L$-embedding 
\begin{equation}\label{equation: restriction of L-groups}
    \xi: \prescript{L}{}{M} \hookrightarrow \prescript{L}{}{G}.
\end{equation} 

%\todo[inline]{$G$ quasisplit has a pure inner form $G_\delta$; this $\delta$ determines a set of $\delta_M$. Write this subtlety somewhere, probably here. See Lemma \ref{lemma: galois cohomology map}\\ \textbf{[disposition:} written, immediately below.\textbf{]}}

Let $\lambda$ be an infinitesimal parameter for $G$ which factors through an infinitesimal parameter $\lambda_M$ of $M$:
\begin{equation}\label{equation: setting inf par factors}% https://q.uiver.app/#q=WzAsMyxbMCwwLCJXX0YiXSxbMiwwLCJcXExncm91cHtHfSJdLFsxLDEsIlxcTGdyb3Vwe019Il0sWzAsMSwiXFxsYW1iZGEiXSxbMCwyLCJcXGxhbWJkYV9NIiwyXSxbMiwxLCJcXHhpIiwyLHsic3R5bGUiOnsidGFpbCI6eyJuYW1lIjoiaG9vayIsInNpZGUiOiJ0b3AifX19XV0=
\begin{tikzcd}
	{W_F} && {\Lgroup{G}} \\
	& {\Lgroup{M}}
	\arrow["\lambda", from=1-1, to=1-3]
	\arrow["{\lambda_M}"', from=1-1, to=2-2]
	\arrow["\xi"', hook, from=2-2, to=1-3]
\end{tikzcd}
\end{equation}
The $L$-embedding \eqref{equation: restriction of L-groups} restricts to an
embedding of the complex algebraic groups:
\begin{equation}\label{equation: embedding of complex dual groups}
    \xi: \dualgroup{M} \hookrightarrow \dualgroup{G}.
\end{equation}

Let $M_\gamma$ (resp. $G_\delta$) correspond to the pure inner form $\gamma \in H^1(F,M)$ (resp. $\delta \in H^1(F,G)$). The inclusion $\xi: \dualgroup{M} \hookrightarrow \dualgroup{G}$ induces a homomorphism $H_\xi: H^1(F,M) \xrightarrow{} H^1(F,G) $, which we realize using Kottwitz isomorphism \eqref{eqn: kottwitz map} below. The inclusion $\xi$ induces the inclusion $Z(\dualgroup{G})^{\Gamma_F} \hookrightarrow Z(\dualgroup{M})^{\Gamma_F}$ which gives the map
    \[\pi_0(Z(\dualgroup{G})^{\Gamma_F}) \rightarrow \pi_0(Z(\dualgroup{M})^{\Gamma_F}).\]
    This in turn induces the map 
    \begin{align}\label{eq: Z xi}
        Z_\xi: \Hom(\pi_0(Z(\dualgroup{M})^{\Gamma_F}), \mathbb{C}^\times) \rightarrow \Hom(\pi_0(Z(\dualgroup{G})^{\Gamma_F}), \mathbb{C}^\times).
    \end{align}
    Then we obtain the following diagram.
  % https://q.uiver.app/#q=WzAsNCxbMCwwLCJIXjEoRixNKSJdLFswLDEsIlxcb3BlcmF0b3JuYW1le0hvbX0oXFxwaV8wKFooXFx3aWRlaGF0e019KSksIFxcbWF0aGJie0N9XlxcdGltZXMpIl0sWzEsMSwiXFxvcGVyYXRvcm5hbWV7SG9tfShcXHBpXzAoWihcXHdpZGVoYXR7R30pKSwgXFxtYXRoYmJ7Q31eXFx0aW1lcykiXSxbMSwwLCJIXjEoRixHKSJdLFsxLDIsIlpfXFx2YXJlcHNpbG9uIiwyXSxbMiwzLCJrX0ciLDJdLFswLDEsImtfTV57LTF9IiwyXSxbMCwzLCJIX1xcdmFyZXBzaWxvbiIsMCx7InN0eWxlIjp7ImJvZHkiOnsibmFtZSI6ImRhc2hlZCJ9fX1dXQ==
\begin{equation}\label{equation: definition of h xi}
\begin{tikzcd}
	{H^1(F,M)} & {H^1(F,G)} \\
	{\operatorname{Hom}(\pi_0(Z(\dualgroup{M})^{\Gamma_F}), \mathbb{C}^\times)} & {\operatorname{Hom}(\pi_0(Z(\dualgroup{G})^{\Gamma_F}), \mathbb{C}^\times)}
	\arrow["{H_\xi}", dashed, from=1-1, to=1-2]
	\arrow["{\kappa_M^{-1}}"', from=1-1, to=2-1]
	\arrow["{Z_\xi}"', from=2-1, to=2-2]
	\arrow["{\kappa_G}"', from=2-2, to=1-2]
\end{tikzcd}
\end{equation}
Finally, it follows from \cite{Kot83}*{Proposition, part (6)} that
\begin{align}\label{eqn Kottwitz sign Levi}
    e(\gamma)=e(H_\xi(\gamma)).
\end{align}


% A class $\delta \in H^1(F,G)$ does not by itself single out a pure inner form of $M$: what it determines is the fibre $H_\xi^{-1}(\delta) \subseteq H^1(F,M)$, and each $\gamma$ in that fibre gives a Levi subgroup $M_\gamma$ of $G_\delta$. Subsequently, we fix a $\gamma \in H^1(F,M)$, which determines $\delta=H_\xi(\gamma)$, and we work with such pairs.
\subsection{Restriction of sheaves}\label{subsec: restriction of sheaves}

Recall from \eqref{equation: h lambda} that $H_\lambda=Z_{\dualgroup{G}}(\lambda)$, so  $\xi$ from \eqref{equation: embedding of complex dual groups} restricts to the inclusion
\begin{equation}\label{equation: inclusion h lambda}
    \xi: H_{\lambda_M} \hookrightarrow H_\lambda.
\end{equation}
On the other hand, the map $\xi$ also induces a map $\operatorname{Lie}(\dualgroup{M})\to \operatorname{Lie}(\dualgroup{G})$ which restricts to an inclusion of Vogan varieties
\begin{equation}\label{equation: inclusion of vogan varieties}
    d\xi: V_{\lambda_M} \hookrightarrow{} V_\lambda.
\end{equation}
This map $d\xi$ is equivariant with respect to the $H_{\lambda_M}$-action on $V_{\lambda_M}$ and $H_\lambda$-action on $V_\lambda$:
\[d\xi(h\cdot v)=\xi(h)\cdot d\xi(v) \hspace{0.3cm} \forall h\in H_{\lambda_M}, v\in V_{\lambda_M}.\]
For any $x \in D \subset V_{\lambda_M}$, where $D$ is an $H_{\lambda_M}$-orbit in $V_{\lambda_M}$, the $H_\lambda$-orbit $C$ of $d\xi(x)\in V_\lambda$ is called the \emph{saturation} of $D$ in $V_\lambda$. From \eqref{equation: inclusion h lambda} we obtain the embedding $Z_{H_{\lambda_M}}(x) \hookrightarrow Z_{H_\lambda}(d\xi(x))$ which induces the map
\begin{equation}\label{equation: map of fundamental groups}
    A_D^C(\xi): A_D \to A_C
\end{equation}
between the equivariant fundamental groups of the orbits $D \subset V_{\lambda_M}$ and $C\subset V_\lambda$. Note the dependence on the choice of basepoint $x\in D$ (which determines $d\xi(x)\in C$). 

For a finite dimensional representation $\tau: A_C \to \GL(V_\tau)$, the pullback $\tau\vert_{A_D}$ is the representation $\tau \circ A_D^C(\xi)$.
\begin{equation}\label{equation: pullback of a rep of fund group}
%https://q.uiver.app/#q=WzAsMyxbMCwwLCJBX0QiXSxbMSwwLCJBX0MiXSxbMSwxLCJHTChWX1xcdGF1KSJdLFswLDEsIkFfRF5DKFxceGkpIl0sWzEsMiwiXFx0YXUiXSxbMCwyLCJcXHRhdVxcdmVydF97QV9EfSIsMl1d
\begin{tikzcd}
	{A_D} & {A_C} \\
	& {\GL(V_\tau)}
	\arrow["{A_D^C(\xi)}", from=1-1, to=1-2]
	\arrow["{\tau\vert_{A_D}}"', from=1-1, to=2-2]
	\arrow["\tau", from=1-2, to=2-2]
\end{tikzcd}
\end{equation}
%\todo[inline]{draw commutation symbol inside the diagram?\\ \textbf{[disposition:} not done --- the diagram is unambiguous as it stands.\textbf{]}}
The map \eqref{equation: inclusion of vogan varieties} also induces the equivariant pullback 
\[\xi^*:D_{H_\lambda}(V_\lambda) \to D_{H_{\lambda_M}}(V_{\lambda_M}) \]
between equivariant derived categories; see \cite{BL94} for the construction of the equivariant derived category and of the pullback functor along an equivariant map. 

% note{ {the key \texttt{BL94} (Bernstein--Lunts, \emph{Equivariant sheaves and functors}) still has to be added to \texttt{references.bib}.}}
%\todo[inline]{add ref for equivariant restriction\\ \textbf{[disposition:} Bernstein--Lunts cited above; key still to add. Note the direction of $\xi^*$ was reversed here.\textbf{]}}
Recall from Section \ref{subsec: grothendieck groups} that the Grothendieck groups coincide: $K\Per_{H_\lambda}(V_\lambda)=KD_{H_\lambda}(V_\lambda)$. This allows us to define the linear map which we also denote as $\xi^*$:
\begin{equation}\label{equation: restriction linear map perverse sheaves}
    \xi^*: K\Per_{H_\lambda}(V_\lambda) \xrightarrow{} K\Per_{H_{\lambda_M}}(V_{\lambda_M}).
\end{equation} 
Recall from the discussion in Section \ref{subsec: grothendieck groups} that the collection of standard sheaves $\stsh{\tau}{C}$ forms a basis of $K\textrm{Sh}_{H_{\lambda}}(V_{\lambda})= K \Per_{H_{\lambda}}(V_{\lambda})$. Therefore, the linear map $\xi^{\ast}$ is determined by its value on these standard sheaves, which is described explicitly in \cite{ABV92}*{Proposition 23.7(b)}. We recall the formula in the following proposition.

\begin{prop}\label{prop: equivariant restriction}
   In the setting above, we have
    \begin{equation}\label{eqn: restriction formula on standards}
    \xi^*(\stsh{\tau}{C}) = \sum_{\xi(D)\subset C,~\sigma \in \irrepdual{A}_D} m(\sigma,\tau) \cdot \stsh{\sigma}{D}.
\end{equation}
Here $D$ runs over the $H_{\lambda_M}$-orbits in $V_{\lambda_M}$ whose image under $d\xi$ is contained in $C$, $\sigma$ runs over $\irrepdual{A}_D$, and $m(\sigma,\tau)$ denotes the multiplicity of $\sigma$ in $\tau\vert_{A_D}$, in the notation of \eqref{equation: pullback of a rep of fund group}.
\end{prop}


\begin{remark}
   If $D=C_{\phi_M}$ and $C=C_\phi$ for Langlands parameters $\phi_M$ of $M$ and $\phi=\xi\circ\phi_M$ of $G$, then under the identifications \eqref{eqn: equivalence of component group and eq fund group} the map $A_D^C(\xi)$ is the map $A_{\phi_M}^\phi(\xi)$ of \eqref{equation: pullback of tau}. As a result, the multiplicities $m(\sigma,\tau)$ above agree with those of Section \ref{subsec: Induction of tempered representations}. 
\end{remark} 
\subsection{Lift}\label{subsection: lift}
For $\delta\in H^1(F,G)$, define the obvious inclusion and projection
\begin{equation}\label{equation: inclusion of grothendieck groups of reps}
    \iota_\delta^G: K\Pi_{\lambda}(G_\delta) \to K\Pi_\lambda^\pure(G), \ \ \pr_\delta^G:K\Pi_\lambda^\pure(G) \to K\Pi_{\lambda}(G_\delta). 
\end{equation}
Note that for any $\mathcal{F}\in K\Per_{H_\lambda}(V_\lambda)^\delta$,
\begin{equation}\label{equation: transpose of projection}
\langle \pr^G_\delta(\pi), \mathcal{F}\rangle = \langle \pi, \mathcal{F}\rangle.
\end{equation}
We often drop the superscript $G$ and write $\iota_\delta$ and $\pr_\delta$ when the group is clear.
With respect to the pairing in \eqref{eqn:pairing}, we have the transpose of pullback: 

\[\prescript{t}{}{\xi^*}: K\Pi_{\lambda_M}^\pure(M) \xrightarrow{} K\Pi^\pure_\lambda(G),\] characterized by the requirement that, for all $\mathcal{F} \in K\Per_{H_{\lambda}}(V_{\lambda})$ and $\pi \in K\Pi_{\lambda_M}^\pure(M) $, 
\begin{equation}\label{equation: transpose of restriction}
    \langle \prescript{t}{}{\xi^*}(\pi), \mathcal{F} \rangle = \langle \pi, \xi^*(\mathcal{F})\rangle.
\end{equation}

Now define
\begin{equation}\label{equation: lift}
 \Lift_{M_\gamma}^{G_\delta}:=\operatorname{pr}^G_\delta \circ \prescript{t}{}{\xi}^* \circ\iota^M_\gamma.
\end{equation}
That is,
% https://q.uiver.app/#q=WzAsNCxbMCwwLCJLXFxQaV97XFxsYW1iZGFfe019fShNKSJdLFsxLDAsIktcXFBpX1xcbGFtYmRhKEcpIl0sWzAsMSwiS1xcUGlfe1xcbGFtYmRhX3tNfX1ee1xccHVyZX0oTSkiXSxbMSwxLCJLXFxQaV5cXHB1cmVfXFxsYW1iZGEoRykiXSxbMCwyLCJcXGlvdGEiLDIseyJzdHlsZSI6eyJ0YWlsIjp7Im5hbWUiOiJob29rIiwic2lkZSI6ImJvdHRvbSJ9fX1dLFsyLDMsIlxccHJlc2NyaXB0e3R9e317XFx4aX0iLDJdLFszLDEsIlxcb3BlcmF0b3JuYW1le3ByXzB9IiwyXSxbMCwxLCJcXExpZnRfTV5HIl1d
\[\begin{tikzcd}
	{K\Pi_{\lambda_{M}}(M_\gamma)} & {K\Pi_\lambda(G_\delta)} \\
	{K\Pi_{\lambda_{M}}^{\pure}(M)} & {K\Pi^\pure_\lambda(G)}
	\arrow["{\Lift_{M_\gamma}^{G_\delta}}", from=1-1, to=1-2]
	\arrow["\iota_\gamma"', hook', from=1-1, to=2-1]
	\arrow["{\prescript{t}{}{\xi^*}}"', from=2-1, to=2-2]
	\arrow["{\operatorname{pr_\delta}}"', from=2-2, to=1-2]
\end{tikzcd}.\]

\subsection{Parabolic induction of representations}\label{subsec: parabolic induction of reps linear map}
Assume we are in the setting \ref{subsection: setting}. Recall from Section  {\ref{subsec: grothendieck groups}} the Grothendieck group $K\Pi_{\lambda}^\pure(G)$ and the parabolic induction functor $\Ind_P^G$ from Section \ref{subsec: parabolic induction}. Assuming Desideratum \ref{desideratum: Haines}, the linear map $\Ind_M^G$ restricts to
\begin{align}\label{equation: induction definition}
    \Ind_M^G: K\Pi^\pure_{\lambda_M}(M) &\xrightarrow{} K\Pi^\pure_{\lambda}(G),\\
    \pi &\mapsto \Ind_M^G(\pi). {\notag}
\end{align}
Here if $\pi \in K\Pi(M_{\gamma})$, then $\Ind_M^G(\pi)$ means $\Ind_{M_{\gamma}}^{G_{H_{\xi}(\gamma)}}(\pi)$ by abuse of notation. 


Suppose $\phi$ is a Langlands parameter of $G$ that factors through the Langlands parameter $\phi_{M}$ of $M$ via $\xi$.
\begin{equation} {\label{equation: setting L par factors}}% https://q.uiver.app/#q=WzAsMyxbMCwwLCJXX0YiXSxbMiwwLCJcXExncm91cHtHfSJdLFsxLDEsIlxcTGdyb3Vwe019Il0sWzAsMSwiXFxsYW1iZGEiXSxbMCwyLCJcXGxhbWJkYV9NIiwyXSxbMiwxLCJcXHhpIiwyLHsic3R5bGUiOnsidGFpbCI6eyJuYW1lIjoiaG9vayIsInNpZGUiOiJ0b3AifX19XV0=
\begin{tikzcd}
	{W_F\times \SL_2(\BC)} && {\Lgroup{G}} \\
	& {\Lgroup{M}}
	\arrow["\phi", from=1-1, to=1-3]
	\arrow["{\phi_M}"', from=1-1, to=2-2]
	\arrow["\xi"', hook, from=2-2, to=1-3]
\end{tikzcd}
\end{equation}
Recall from Section \ref{subsec: Induction of tempered representations} the homomorphism of component groups $A_{\phi_M}^\phi(\xi):A_{\phi_M}\to A_\phi$ induced by $\xi$, the pullback $\tau\vert_{A_{\phi_M}}:=\tau\circ A_{\phi_M}^\phi(\xi)$ of $\tau \in \irrepdual{A}_\phi$, and the multiplicities $m(\tau_M,\tau)$ of \eqref{equation: pullback of tau}.
%; nothing there required $\phi_M$ to be tempered. 
Recall also the notation $\kappa_G(\tau)$ for $\tau \in \irrepdual{A}_\phi$ from \eqref{eqn: kappa tau} defined using the Kottwitz isomorphism and the map $H_\xi: H^1(F,M)\to H^1(F,G)$ from \eqref{equation: definition of h xi}. 
%\todo[inline]{check if assumptions are written correctly\\ \textbf{[disposition:} reviewed; statement kept as is.\textbf{]}}
\begin{lemma}\label{lemma: relating galois colomology class of irreps of component groups}
     Suppose the irreducible representation $\tau_M\in \irrepdual{A}_{\phi_M}$ occurs in the pullback $\tau\vert_{A_{\phi_M}}$ according to \eqref{equation: pullback of tau}. Then, $\kappa_G(\tau)=H_\xi(\kappa_M(\tau_M))$.
\end{lemma}
\begin{proof}
    For $\tau \in \irrepdual{A}_\phi$, let $\chi_\tau \in \Hom(\pi_0(Z(\dualgroup{G})^{\Gamma_F}), \CC^\times)$ be the character in the pullback of $\tau$ via the map in $\iota_\phi$ in \eqref{equation: inclusion of componenet grp of center in a phi}. By definition \eqref{eqn: kappa tau}, $\kappa_G(\tau)=\kappa_G(\chi_\tau)$. Similarly $\kappa_M(\tau_M)=\kappa_M(\chi_{\tau_M})$. The conclusion follows from the diagram \eqref{equation: definition of h xi}.  
\end{proof}



%\todo[inline]{check if this theorem and proof are written in the intended way\\ \textbf{[disposition:} several corrections made in the statement and the proof --- see the change log (misplaced parentheses in the induction-in-stages steps, $m(\tau,\tau_M)\to m(\tau_M,\tau)$, missing argument of $A^{\phi_M}_{\phi_{M_0}}$, coroots/unipotent radical, ${}^{\widehat w\tau}\to{}^{\widehat w}\tau$).\textbf{]}}


% \todo[inline]{Chi-Heng:This is too much description. I will remove this paragraph and rewrite the statement of the lemma.}

% It allows a \rev{Langlands datum} \todo[inline]{has this been defined somewhere or is it very common terminology?} for $M$ to be moved into standard position inside $G$; the point is that a Weyl conjugate of a standard Levi subgroup need not be standard. Write $A_0$ for the maximal split subtorus of the maximal torus $T$ fixed in Section \ref{Groups and representations}, $\Delta_0$ for the set of simple roots of $A_0$ in $G$, so that the standard parabolic subgroups of $G$ correspond to the subsets of $\Delta_0$, and $W(G,A_0)=N_G(A_0)(F)/Z_G(A_0)(F)$ for the relative Weyl group, every element of which admits a representative in $G(F)$. For a standard Levi subgroup $M_0$ with maximal split central torus $A_{M_0}$ we identify, as usual, the unramified characters occurring in the Langlands classification with the points of $\mfr{a}^*_{M_0}=X^*(A_{M_0})\otimes \R$. 
%positivity referring to the pairings $\langle\,\cdot\,,\alpha^\vee\rangle$ with $\alpha\in\Delta_0$.}
% note{ {the symbol $\Delta$ is doing double duty in this paper: the set of simple roots in Section \ref{Groups and representations} and the standard module $\Delta(\pi)$ everywhere else. I have written $\Delta_0$ for the relative simple roots here, but you may want to rename one of the two.}}

Here is the main result of this section.

\begin{thm}\label{theorem: induction formula}
Assume that there exists an $\LLC^{\mathfrak{e}}$ for Levi subgroups of $G$ satisfying Desiderata \ref{desideratum: Langlands classification}, \ref{desideratum: induction of tempered representations}, \ref{desideratum: isomorphism}, and \ref{desideratum: Haines}. Recall the notation set up in Section \ref{subsec: Induction of tempered representations}. Suppose now that $\phi$ is \emph{any} $L$-parameter that factors through the parameter $\phi_M$ of a standard Levi subgroup $M$ of $G$. We have in the Grothendieck group $K\Pi_\lambda^\pure(G)$
\begin{equation}\label{induction formula}
   \Ind_{M}^{G} (\Delta(\phi_M,\tau_M))=\sum_{\tau \in \widehat{A}_{\phi}} m(\tau_M, \tau) \Delta(\phi, \tau).
\end{equation}
\end{thm}
\begin{proof}
    Write $\phi_M \leftrightarrow (P_0, \phi_{M_0}, \nu_0)$ as an $L$-parameter of $M$. Thus $P_0$ is a standard parabolic subgroup of $M$ with standard Levi component $M_0$, the parameter $\phi_{M_0}$ is tempered, and $\nu_0 \in \mfr{a}^*_{M_0}$ is positive with respect to $P_0$.  By Proposition \ref{prop: SZ18}, we can assume
    %now we regard $\phi_M$ as a  Langlands parameter such that 
    ${}^LM_0 = Z_{{}^L M}(z_{\phi_M})$, where $z_{\phi_M}$ is the hyperbolic part of $\phi_M(\Fr) \in {}^L{M}$ for any Frobenius element $\Fr$ of $F$. Let $\xi_{M_0}^M: \prescript{L}{}{M_0} \hookrightarrow \prescript{L}{}{M}$ be the $L$-embedding, and let $\tau_{M_0}:= \tau_M \circ A_{\phi_{M_0}}^{\phi_M}(\xi_{M_0}^M)$, which is irreducible by Proposition \ref{prop: SZ18}(3). In particular, $\tau_{M_0} \in \irrepdual{A}_{\phi_{M_0}}$. Let $P$ be the standard parabolic subgroup of $G$ with $P\cap M=P_0$, so that $M_0$ is the Levi component of $P$. Inducing in stages, the left hand side of \eqref{induction formula} is
    \[ \Ind_{M}^G \big(\Ind_{M_0}^M (\pi(\phi_{M_0},\tau_{M_0}) \otimes \nu_0)\big) = \Ind_{M_0}^G (\pi(\phi_{M_0},\tau_{M_0}) \otimes \nu_0).\]
    We show that the right-hand side is a sum of standard modules (with multiplicities) in the Grothendieck group.

    Although $\nu_0$ is positive with respect to $P_0$, it is not necessarily positive with respect to $P$. 
    Therefore, we take an element $w_1 \in N_G(A_0)(F)$, where $A_0$ is the maximal split torus of the maximal torus $T$ fixed in Section \ref{Groups and representations}, such that $w_1(\nu_0)$ lies in the closure of the Weyl chamber determined by the Borel subgroup we fixed. Then the character $w_1(\nu_0)$ of $w_1M_0w_1^{-1}$ extends to a character of a standard Levi subgroup $M_1:= Z_{G}(w_1(\nu_0))$, which is positive with respect to the corresponding standard parabolic subgroup $P_1$. Since $w_1M_0w_1^{-1}$ contains the centralizer of $A_0$, it is a semi-standard Levi subgroup of $M_1$. We may take $w_2 \in N_{M_1}(A_0)$ such that $M_0':= w_2w_1M_0w_1^{-1}w_2^{-1}$ is a standard Levi subgroup of $M_1$ (and of $G$). Set $w:=w_2w_1$. For a representation $\sigma$ of a pure inner form $M_{0,\gamma_0}$ of $M_0$, we let $\sigma^w$ denote the corresponding representation of the pure inner form $M_{0,w_{\ast}\gamma_0}'$ of $M_0'$. Note that the image of $\gamma_0 \in H^1(F,M_0) $ and $w_{\ast}\gamma_0 \in H^1(F,M_0')$ in  $ H^1(F,G)$ agree since $w \in G(F)$.
    Now we can continue to rewrite the left hand side of \eqref{induction formula} in Grothendieck group as  
\begin{align*}
     \Ind_{M_0}^G (\pi(\phi_{M_0},\tau_{M_0}) \otimes \nu_0)&=\Ind_{w M_0 w^{-1}}^G (\pi(\phi_{M_0},\tau_{M_0})^w \otimes w(\nu_0))\\
     &=\Ind_{M_1}^G (\Ind_{M_0'}^{M_1} (\pi(\phi_{M_0},\tau_{M_0})^w \otimes w(\nu_0)))\\
     &=\Ind_{M_1}^G (\Ind_{M_0'}^{M_1} (\pi({}^{\widehat{w}}\phi_{M_0},{}^{\widehat{w}}\tau_{M_0}) )\otimes w(\nu_0)).
\end{align*}
For the last equality, we use Desideratum \ref{desideratum: isomorphism} and the notation discussed below it. Now we decompose the tempered induction $\Ind_{M_0'}^{M_1}\pi({}^{\widehat{w}}\phi_{M_0},{}^{\widehat{w}}\tau_{M_0})$ using Desideratum \ref{desideratum: induction of tempered representations} explicitly. In the following, if $Q_1 \leq Q_2 \leq G$ are standard Levi subgroups, we let $\xi_{Q_1}^{Q_2}$ denote the $L$-embedding from ${}^L Q_1$ to ${}^L Q_2$.    Set $\phi'={\xi_{M_0'}^{M_1}} \circ {}^{\widehat{w}} \phi_{M_0}$. We have
\[\Ind_{M_0'}^{M_1}\pi({}^{\widehat{w}}\phi_{M_0},{}^{\widehat{w}}\tau_{M_0})= \bigoplus_{\tau' \in \widehat{A}_{\phi'} } m( {}^{\widehat{w}}\tau_{M_0},\tau') \pi(\phi', \tau').\]
Each $\pi(\phi',\tau')$ is a tempered representation of $M_1$ and $w(\nu_0)$ is positive with respect to $P_1$, so  $\Ind_{P_1}^G(\pi(\phi',\tau')\otimes w(\nu_0))$ is a standard representation. Put
\[\phi'':= \xi_{M_1}^G \circ \big(\phi' \otimes w(\nu_0)\big) = \xi_{M_0'}^G \circ \big({}^{\widehat{w}} \phi_{M_0}\otimes w(\nu_0)\big).\]
By Desideratum \ref{desideratum: Langlands classification} and Proposition \ref{prop: SZ18}(3) that standard module is $\Delta(\phi'',\tau'')$, where $\tau''\in\irrepdual{A}_{\phi''}$ corresponds to $\tau'$ under the isomorphism $A_{\phi'}\simeq A_{\phi''}$ induced by $\xi_{M_1}^G$. As $\xi_{M_0'}^G=\xi_{M_1}^G\circ \xi_{M_0'}^{M_1}$, this isomorphism is compatible with the two maps from $A_{{}^{\widehat{w}}\phi_{M_0}}$, so that $m({}^{\widehat{w}}\tau_{M_0},\tau')=m({}^{\widehat{w}}\tau_{M_0},\tau'')$. We obtain in the Grothendieck group that
\begin{align*}
    \Ind_M^G \Delta(\phi_M, \tau_M)&= \sum_{\tau' \in \widehat{A}_{\phi'} }   m( {}^{\widehat{w}}\tau_{M_0},\tau')  \Ind_{{M_1}}^G \big( \pi(\phi', \tau') \otimes w(\nu_0)\big)\\
    &=\sum_{\tau'' \in \widehat{A}_{\phi''} }   m( {}^{\widehat{w}}\tau_{M_0},\tau'')  \Delta(\phi'', \tau'').
\end{align*}

It remains to compare $\phi''$ with $\phi$. The Langlands classification of $\phi_M$ reads $\phi_M=\xi_{M_0}^M\circ(\phi_{M_0}\otimes\nu_0)$, so that $\phi=\xi_M^G\circ \phi_M=\xi_{M_0}^G\circ(\phi_{M_0}\otimes\nu_0)$. Comparing with the second expression for $\phi''$ above, we see that $\phi''={}^{\widehat{w}}\phi$. Hence $\phi''=\phi$ as $L$-parameters of $G$, and the isomorphism $A_{\phi''}\to A_\phi$ is conjugation by $\widehat{w}^{-1}$. Thus, we may rewrite
\begin{align*}
      \Ind_M^G \Delta(\phi_M, \tau_M)
    &=\sum_{\tau'' \in \widehat{A}_{\phi''} }   m( {}^{\widehat{w}}\tau_{M_0},\tau'')  \Delta(\phi'', \tau'')\\
    &=\sum_{\tau \in \widehat{A}_{\phi} }   m( {}^{\widehat{w}}\tau_{M_0},{}^{\widehat{w}}\tau)  \Delta(\phi, \tau)\\
    &=\sum_{\tau \in \widehat{A}_{\phi} }   m( \tau_{M_0},\tau)  \Delta(\phi, \tau)\\
    &=\sum_{\tau \in \widehat{A}_{\phi} }   m( \tau_{M},\tau)  \Delta(\phi, \tau).
\end{align*}
Here, the fourth equality comes from the isomorphism $A_{\phi_{M_0}}\simeq A_{\phi_M}$ by Proposition \ref{prop: SZ18}(3), which identifies $\tau_{M_0} \in \widehat{A}_{\phi_{M_0}}$ with $\tau_M \in \widehat{A}_{\phi_M}$. This proves \eqref{induction formula}.

{Finally we deduce the refinement. If $\kappa_M(\tau_M)=\gamma$ and $m(\tau_M,\tau)\neq 0$, then $\kappa_G(\tau)=H_\xi(\gamma)$ by Lemma \ref{lemma: relating galois colomology class of irreps of component groups}; so only such $\tau$ contribute to \eqref{induction formula}, which is \eqref{equation: induction formula form by form}, and the map $\Ind_{M_\gamma}^{G_\delta}$ is well defined. This completes the proof.}
\end{proof}


\begin{remark}
Though \eqref{induction formula} is an equation in $K\Pi_{\lambda}^{\pure}(G)$, the representations involved are in the same pure inner form of $G$. More precisely, let $\gamma:= \kappa_M(\tau_M)\in H^1(F,M)$ and let $\delta:= H_{\xi}(\gamma) \in H^1(F,G)$. By Lemma \ref{lemma: relating galois colomology class of irreps of component groups}, if $m(\tau_M, \tau) \neq 0$, then $\kappa_G(\tau)= \delta$. In conclusion, we may rewrite \eqref{induction formula} form-by-form:
    \begin{equation}\label{equation: induction formula form by form}
        \Ind_{M_\gamma}^{G_\delta}(\Delta(\phi_M,\tau_M))=\sum_{\tau \in \irrepdual{A}_\phi,~\kappa_G(\tau)=H_\xi(\gamma)}{m(\tau_M, \tau)} \Delta(\phi, \tau).
    \end{equation}
\end{remark}


\section{Parabolic induction as the transpose of geometric restriction}
Throughout this section, we are in the setting of \S \ref{subsection: setting}. In Section \ref{subsection: pure parabolic induction and restriction}  we show that $\Ind_M^G=\prescript{t}{}{\xi^*}$, where the transpose is taken with respect to the pairing \eqref{eqn:pairing} determined by $\LLC^{\mathfrak{e}}$. For $\gamma\in H^1(F,M)$ and $\delta=H_\xi(\gamma)$, we also offer the obvious refinement of this identity in Section \ref{subsection: parabolic induction form-by-form and lift}, by showing that $\Ind_{M_\gamma}^{G_\delta}$ is $\Lift_{M_\gamma}^{G_\delta}$, where the latter is defined in Section \ref{subsection: lift}.

\subsection{Pure parabolic induction and restriction}\label{subsection: pure parabolic induction and restriction}

%\section{The duality of parabolic induction and restriction of perverse sheaves}

\begin{prop}\label{proposition: pure parabolic induction and restriction}
    Let $G$ be a quasisplit connected reductive algebraic group with Levi subgroup $M$. We retain the notation from Section \ref{subsection: setting}. Assume that there exists an $\LLC^{\mathfrak{e}}$ for Levi subgroups of $G$ satisfying Desiderata \ref{desideratum: Langlands classification}, \ref{desideratum: induction of tempered representations}, \ref{desideratum: isomorphism}, and \ref{desideratum: Haines}. Further assume the $p$-adic Kazhdan-Lusztig Hypothesis \ref{hypothesis: KLH} for both $G$ and $M$. For $\pi\in K\Pi_{\lambda_M}^\pure(M)$ and $\mathcal{F}\in K\Per_{H_\lambda}(V_\lambda)$, we have
    \begin{equation}\label{equation: ind res adjunction}\langle \Ind_M^G(\pi),\mathcal{F}\rangle = \langle \pi,\xi^*(\mathcal{F})\rangle {.}\end{equation}
{This is the implication \textup{(a)}$\Rightarrow$\textup{(b)} of Theorem \ref{theorem: main theorem in the introduction}.}
\end{prop}
\begin{proof}
    It suffices to check the equality on the bases of standard representations of $K\Pi_{\lambda_M}^\pure(M)$ and of standard constructible sheaves of $K\Per_{H_\lambda}(V_\lambda)$, as provided by \eqref{equation: gr grp perverse sheaves} and \eqref{equation: gr grp of reps}. Consider the standard representation $\Delta(\phi_M,\tau_M)$ attached to the enhanced $L$-parameter $(\phi_M,\tau_M)$ of $M$, where $\tau_M\in \irrepdual{A}_{\phi_{M}}$. Let $\mathcal{F}=\stsh{\tau}{C}$, where $C$ is an $H_\lambda$-orbit in $V_\lambda$ and $\tau \in \irrepdual{A}_C$. From Theorem \ref{theorem: induction formula}, we have that 
    \[\Ind_{M}^{G} (\Delta(\phi_M,\tau_M))=\sum_{\tau' \in \widehat{A}_{\phi}} m(\tau_M, \tau') \Delta(\phi, \tau').\]
    This, combined with the identification $A_\phi=A_{C_\phi}$ from \eqref{eqn: equivalence of component group and eq fund group}, where $C_\phi$ is the $H_\lambda$-orbit in $V_\lambda$ corresponding to $\phi$, we get  
    \begin{align*}
    {\langle\Ind_M^G(\Delta(\phi_M,\tau_M)),\stsh{\tau}{C}\rangle}
        &=   \displaystyle \sum_{\tau' \in \widehat{A}_{C_\phi}} m(\tau_M, \tau')\langle \Delta(\phi, \tau'), \stsh{\tau}{C}\rangle \\
          &=\begin{cases}
             m(\tau_M,\tau)\langle \Delta(\phi,\tau),\stsh{\tau}{C} \rangle & \text{ if } C=C_\phi,\\
             0 & \text{otherwise}.
          \end{cases}  \\
          &=\begin{cases}
             m(\tau_M,\tau) e(\kappa_G(\tau)) & \text{ if } C=C_\phi,\\
             0 & \text{otherwise}. 
          \end{cases}  
    \end{align*}
where the last equality is evaluated using the Kazhdan-Lusztig Hypothesis (Hypothesis \ref{hypothesis: KLH}). This computes the left hand side of \eqref{equation: ind res adjunction}.

Now we compute the right hand side of \eqref{equation: ind res adjunction}.
From  \eqref{eqn: restriction formula on standards}, we have
\[\xi^*(\stsh{\tau}{C}) = \sum_{\xi(D)\subset C,~\sigma \in \irrepdual{A}_D} m(\sigma,\tau) \cdot \stsh{\sigma}{D}.\] 
We obtain
\begin{align*}
     \langle \Delta(\phi_M,\tau_M), {\xi^*(\stsh{\tau}{C})}\rangle &= \displaystyle\sum_{\xi(D)\subset C,~\sigma \in \irrepdual{A}_D}m(\sigma,\tau)\langle \Delta(\phi_M,\tau_M),\stsh{\sigma}{D}\rangle \\&=\begin{cases}
         m(\tau_M,\tau) e(\kappa_M(\tau_M)) & \text{ if } C=C_\phi,\\
         0 & \text{ otherwise.}
         \end{cases}
\end{align*}
Here the second equality holds because a summand $\langle \Delta(\phi_M,\tau_M),\stsh{\sigma}{D}\rangle$ is nonzero only for $(D,\sigma)=(C_{\phi_M},\tau_M)$, which forces $C=C_\phi$, the saturation of the orbit $C_{\phi_M}$ in $V_\lambda$. The resulting value is then read off once again using Hypothesis \ref{hypothesis: KLH}. If $m(\tau_M,\tau)\neq 0$, then $\kappa_G(\tau)=H_\xi(\kappa_M(\tau_M))$ by Lemma \ref{lemma: relating galois colomology class of irreps of component groups}, whence $e(\kappa_G(\tau))=e(\kappa_M(\tau_M))$ by \eqref{eqn Kottwitz sign Levi}. 
\end{proof}


\subsection{Parabolic induction form-by-form and lift}\label{subsection: parabolic induction form-by-form and lift}


Recall the map \[\Lift_{M_\gamma}^{G_\delta}: K\Pi_{\lambda_M}(M_\gamma)\to K\Pi_\lambda(G_\delta)\] from Section \ref{subsection: lift}.
\begin{prop}\label{lemma: lift is induction}
    Assume that there exists a $\LLC^{\mathfrak{e}}$ for Levi subgroups of $G$ satisfying Desiderata \ref{desideratum: Langlands classification}, \ref{desideratum: induction of tempered representations}, \ref{desideratum: isomorphism}, and \ref{desideratum: Haines}. Further assume the $p$-adic Kazhdan-Lusztig Hypothesis \ref{hypothesis: KLH} both for $G$ and $M$. We retain the notation from Sections \ref{subsection: setting} and \ref{subsection: lift}. We have
    \[ \Ind_{M_\gamma}^{G_\delta} = \Lift_{M_\gamma}^{G_\delta}.\]
\end{prop}
\begin{proof}
    It suffices to show that 
    \begin{equation}\label{eqn: need to show}
        \langle \Ind_{M_\gamma}^{G_\delta}(\Delta(\phi_M, \tau_M)), \stsh{\tau}{C} \rangle = \langle \Lift_{M_\gamma}^{G_\delta}(\Delta(\phi_M, \tau_M)), \stsh{\tau}{C}\rangle
    \end{equation}
    for all $ (\phi_M,\tau_M) \in \Phi^{\mfr{e}}(M)$ such that $\kappa_M(\tau_M)=\gamma$, and all orbits $C \in V_{\lambda}/H_{\lambda}$ and $\tau \in \irrepdual{A}_{C}$.
    
    
    For the left hand side of \eqref{eqn: need to show}, using  \eqref{equation: induction formula form by form}, we have 
    \begin{align*}
        \langle \Ind_{M_\gamma}^{G_\delta}(\Delta(\phi_M, \tau_M)),\stsh{\tau}{C} \rangle &= \sum_{\tau' \in \irrepdual{A}_\phi,~\kappa_G(\tau')=\delta} m(\tau_M,\tau')\langle \Delta(\phi, \tau'), \stsh{\tau}{C} \rangle \\
        &= \begin{cases}
            m(\tau_M,\tau)\langle \Delta(\phi, {\tau}), \stsh{\tau}{C_\phi} \rangle  & \text{if } C=C_\phi \text{ and } \kappa_G(\tau)=\delta,\\
            0 & \text{otherwise}.
        \end{cases}\\
        &= \begin{cases}
            m(\tau_M,\tau)e_G(\delta)  & \text{if } C=C_\phi \text{ and } \kappa_G(\tau)=\delta,\\
            0 & \text{otherwise}.
        \end{cases}
    \end{align*}
The computation of $\langle \Delta(\phi, {\tau}), \stsh{\tau}{C_\phi} \rangle$ in {the} last equality follows from the Kazhdan-Lusztig Hypothesis \ref{hypothesis: KLH}. 

For the right hand side of \eqref{eqn: need to show}, using \eqref{equation: lift}, we have
\begin{align*}
   \langle \Lift_{M_\gamma}^{G_\delta}&(\Delta(\phi_M, \tau_M)), \stsh{\tau}{C}\rangle\\
   &= \big\langle \pr^G_\delta \circ \prescript{t}{}{\xi^*} \circ \iota^M_\gamma (\Delta(\phi_M, \tau_M)), \stsh{\tau}{C}\big\rangle\\[2pt]
   &= \big\langle \pr^G_\delta \circ \prescript{t}{}{\xi^*} (\Delta(\phi_M,\tau_M)),\stsh{\tau}{C}\big\rangle
      && \text{by }\eqref{equation: inclusion of grothendieck groups of reps}\\[2pt]
   &= \begin{cases}
        \big\langle \prescript{t}{}{\xi^*} (\Delta(\phi_M,\tau_M)), \stsh{\tau}{C}\big\rangle & \text{if }\kappa(\tau)=\delta,\\
        0 & \text{otherwise},
      \end{cases}
      && \text{by }\eqref{equation: transpose of projection}\\[2pt]
   &= \begin{cases}
        \big\langle \Delta(\phi_M,\tau_M), \xi^*(\stsh{\tau}{C})\big\rangle & \text{if }\kappa(\tau)=\delta,\\
        0 & \text{otherwise},
      \end{cases}
      && \text{by }\eqref{equation: transpose of restriction}\\[2pt]
   &= \begin{cases}
        m(\tau_M,\tau)\, e_M(\gamma) & \text{if } C=C_\phi \text{ and } \kappa_G(\tau)=\delta,\\
        0 & \text{otherwise},
      \end{cases}
      && \text{as in Proposition \ref{proposition: pure parabolic induction and restriction}.}
\end{align*}
Finally, note that $e_M(\gamma)= e_{G}(H_{\xi}(\gamma))= e_G(\delta)$ by \eqref{eqn Kottwitz sign Levi}. We conclude that \eqref{eqn: need to show} holds. 
\end{proof}

\section{The converse implication: recovering the two conjectures}\label{section: converse}
In this section, we prove the implication \textup{(b)}$\Rightarrow$\textup{(a)} of Theorem \ref{theorem: main theorem in the introduction}. It is the converse of Proposition \ref{proposition: pure parabolic induction and restriction} in the following sense. Recall the maps $\Ind_M^G:K\Pi_{\lambda_M}^\pure(M)\to K\Pi_\lambda^\pure(G)$ from \eqref{equation: induction definition} and $\xi^*:K\Per_{H_\lambda}(V_\lambda)\to K\Per_{H_{\lambda_M}}(V_{\lambda_M})$ from \eqref{equation: restriction linear map perverse sheaves}. Assuming the identity $\Ind_M^G= {}^t \xi^{\ast}$,  we recover both the $p$-adic analogue of the Kazhdan-Lusztig hypothesis and the decomposition of tempered induction in
Desideratum \ref{desideratum: induction of tempered representations}. 

\begin{prop}\label{theorem: ind res adjunction implies klh}
Let $G$ be a quasisplit connected reductive algebraic group.  Assume that there exists an $\LLC^{\mathfrak{e}}$ for Levi subgroups of $G$ satisfying Desiderata \ref{desideratum: Langlands classification} and \ref{desideratum: Haines}. Suppose further that for any standard Levi subgroup $M$ and infinitesimal parameter $\lambda_M$, the identity \eqref{equation: ind res adjunction}
\[\langle \Ind_M^G(\pi),\mathcal{F}\rangle=\langle \pi, \xi^*(\mathcal{F})\rangle  \]
holds for all $\pi \in K\Pi_{\lambda_M}^{\pure}(M)$ and $\mathcal{F} \in K \Per_{H_{\lambda}}(V_{\lambda
})$. Then
\begin{enumerate}
    \item [(a)] The $p$-adic Kazhdan-Lusztig hypothesis holds$:$ 
\begin{equation*}
    \begin{array}{rcl}
        K\Pi_\lambda^\pure(G) \times K\Per_{H_\lambda}(V_\lambda)
            &\to& \mathbb{Z}\\
        \langle \Delta(\pi), \tau_{C_{\phi}}^{} \rangle 
            &\mapsto& 
            \begin{cases}
                e(\pi) & \LLC^{\mathfrak{e}}_G(\pi) = (\phi,\tau), \\
                0 & \text{else}.
            \end{cases}
    \end{array}
\end{equation*}
    \item [(b)]  Desideratum \ref{desideratum: induction of tempered representations} holds: for any standard Levi subgroup $M$ of $G$, and tempered $L$-parameter $\phi_M$ of $M$,
    \[\Ind_{M}^{G} \big(\pi(\phi_M,\tau_M)\big)=\bigoplus_{\tau \in \irrepdual{A}_{\phi}} m(\tau_M, \tau)\, \pi(\phi, \tau).\]
\end{enumerate}
\end{prop}
\begin{proof}
We first prove Part (a).   For an irreducible representation $\pi \in \Pi_{\lambda}^\pure(G)$, write 
    \[\Delta(\pi)=\Ind_M^G(\pi_M) \in K\Pi_{\lambda}^\pure(G),\]
    where $\pi_M$ is an essentially tempered representation in $\Pi_{\lambda_M}^\pure(M)$ with enhanced Langlands parameter $(\phi_M, \sigma)$.  Let $\LLC_G^{\mathfrak{e}}(\pi)=(\phi,\tau)$; then $\phi=\xi\circ\phi_M$ as in \eqref{equation: setting L par factors}, and in particular $\phi$ has infinitesimal parameter $\lambda= \xi \circ \lambda_M$. Moreover, we know from Proposition \ref{prop: SZ18}(3) that 
\begin{equation}\label{equation: isom of comp groups tempered}
    A_\phi \simeq A_{\phi_M}, 
\end{equation}
so we may write \begin{equation}\label{equation: sigma is tau}
  \tau \circ   A_{\phi_M}^{\phi}=\sigma, \text{ or with mild abuse of notation, } \tau=\sigma.
\end{equation}
Let $C_{\phi_M}$ denote the $H_{\lambda_M}$-orbit in $V_{\lambda_M}$ corresponding to $\phi_M$ according to \eqref{eqn: vogan langlands correspondence}. The orbit $C_{\phi_M}$ is open in $V_{\lambda_M}$; see Remark \ref{rmk LC tempered}. Let $C$ be an $H_\lambda$-orbit in $V_\lambda$ and $\gamma \in \irrepdual{A}_C$. By assumption of identity \eqref{equation: ind res adjunction},
    \begin{equation}\label{equation: application of ind res adj}
        \langle \Ind_M^G(\pi_M), \stsh{\gamma}{C}\rangle = \langle \pi_M, \xi^*(\stsh{\gamma}{C})\rangle. \end{equation}
%\todo[inline]{explain this more\\ \textbf{[disposition:} expanded into the argument that follows; the previous ``$C_{\phi_M}$ contained in $\xi^*(C)$'' was not well formed.\textbf{]}}
By \eqref{eqn: restriction formula on standards}, the sheaf $\xi^*(\stsh{\gamma}{C})$ is supported on those $H_{\lambda_M}$-orbits $D$ with $d\xi(D)\subseteq C$. Since the pairing on the right-hand side of \eqref{equation: application of ind res adj} only sees the orbit $C_{\phi_M}$ attached to $\pi_M$, it vanishes unless $d\xi(C_{\phi_M})\subseteq C$. %as $C_{\phi_M}$ is open in $V_{\lambda_M}$ and 
The orbit $C_\phi$ is by definition the saturation of $C_{\phi_M}$ in $V_\lambda$ so $C=C_\phi$.
Recall from \eqref{eqn: equivalence of component group and eq fund group} that $A_{\phi}$ (resp. $A_{\phi_M}$) is the equivariant fundamental group of $C_\phi$ (resp. $C_{\phi_M}$). Using the formula \eqref{eqn: restriction formula on standards} for $\xi^*$ with the identification \eqref{equation: isom of comp groups tempered} in mind, we get
\[\langle \pi_M, \xi^*(\stsh{\gamma}{C})\rangle= \langle \pi_M,\stsh{\gamma}{C_{\phi_M}} \rangle.\]
Let $j:C_{\phi_M}\hookrightarrow V_{\lambda_M}$ denote the inclusion. Since $C_{\phi_M}$ is open in $V_{\lambda_M}$ ($\ol{C}_{\phi_M}=V_{\lambda_M}$), and $\locsys{\gamma}{C_{\phi_M}}[\dim C_{\phi_M}]\in \Per(C_{\phi_M})$, we have from \eqref{equation: standard sheaf as a sum of perverse sheaf and boundary terms} 
\[(-1)^{\dim C_{\phi_M}}\stsh{\gamma}{C_{\phi_M}}=\mathcal{IC}(\locsys{\gamma}{C_{\phi_M}})+ \text{ boundary terms} ~~\in K\Per_{H_{\lambda_M}}(V_{\lambda_M}).\]
%\todo[inline]{omg nontrivial fact, keep an eye in revisions; maybe consider adding in preliminaries in geometry; $j$ need not be affine but it is ok. I don't think boundary terms contribute anything to the pairing but check again since we know all the simple objects.}

Substituting $\stsh{\gamma}{C_{\phi_M}}$ in the pairing, the boundary terms pair trivially with $\pi_M$, so we get

%the perverse extension $\mathcal{IC}(\locsys{\gamma}{C_{\phi_M}})$ of the local system $\locsys{\gamma}{C_{\phi_M}}$ is the shifted extension by zero $\stsh{\gamma}{C_{\phi_M}}[\dim C_{\phi_M}]$ of $\locsys{\gamma}{C_{\phi_M}}$ minus some boundary terms (as an element of the Grothendieck group $K\Pi_{\lambda_M}^\pure(M)$); see \cite[Section 4]{BBD82}. \todo[inline]{this turned out to be subtle and nontrivial; keep an eye}
\begin{align}\label{eq pairing converse}
\nonumber \langle \pi_M, \xi^*(\stsh{\gamma}{C_\phi})\rangle 
&= \left\langle \pi_M, \stsh{\gamma}{C_{\phi_M}}\right\rangle \\
&= (-1)^{\dim C_{\phi_M}}\left\langle \pi_M, \mathcal{IC}\left(\locsys{\gamma}{C_{\phi_M}}\right)\right\rangle \\
 \nonumber &=\begin{cases}
    (-1)^{2\dim C_{\phi_M}}~e(\pi_M) & \text{ if } C=C_\phi \text{ and } \gamma=\sigma,\\
    0 & \text{otherwise,} 
\end{cases}\\
\nonumber   &= \begin{cases}
    e(\pi) & \text{ if } C=C_\phi \text{ and } \gamma=\tau,\\
    0 & \text{otherwise.}
    \end{cases}
\end{align}
where the second to last equality follows from the definition of the pairing \eqref{eqn:pairing}, and in the last equality the condition follows from \eqref{equation: sigma is tau} and $e(\pi_M)=e(\pi)$  from \eqref{eqn Kottwitz sign Levi}.

In summary, we have computed that
\[ \langle \Delta(\pi), \stsh{\gamma}{C} \rangle=\begin{cases}
    e(\pi) & \text{ if } C=C_\phi \text{ and } \gamma=\tau,\\
    0 & \text{otherwise.}
    \end{cases} \]
In the nonzero instance, the condition $\langle \Delta(\pi),{\stsh{\gamma}{C}}\rangle = e(\pi) $ if and only if $ (C,\gamma)=(C_\phi,\tau) $ can be rewritten as $\langle \Delta(\pi),\stsh{\tau}{C_\phi}\rangle = e(\pi) $ if and only if $ \LLC_G^{\mathfrak{e}}(\pi)=(\phi, \tau)$. This completes the proof of Part (a).


Now we prove Part (b). From now on, $\phi_M$ is a tempered $L$-parameter of $M$ with infinitesimal parameter $\lambda_M$, and $\phi=\xi \circ \phi_M, \lambda= \xi \circ \lambda_M$ are the corresponding tempered $L$-parameter and infinitesimal parameter of $G$. It is well known that induction of a tempered representation is a semisimple direct sum of tempered representations. Moreover, by Desideratum \ref{desideratum: Haines}, we know that 
\[ \Ind_{M}^G (\pi(\phi_M,\tau_M))= \bigoplus_{\tau \in \irrepdual{A}_{\phi}} m_{\tau} \pi(\phi,\tau),\]
where $m_{\tau}$ are non-negative integers. Indeed, if $\pi=\pi(\phi',\tau')$ appears in $\Ind_M^G(\pi(\phi_M,\tau_M))$, then 
\[ \lambda_{\phi'}= \xi \circ \lambda_{\phi_M}= \lambda_{\phi}=:\lambda.\]
Since $\phi'$ and $\phi$ are both tempered, they correspond to the unique open orbit in $V_{\lambda}$ (\cite{cunningham2024whittakernormalizationpadicabvpackets}*{Proposition 3.2}). Hence, $\phi'$ is equal to $\phi$. 

It thus remains to show that the multiplicity $m_{\tau}$ matches $ m(\tau_M,\tau)$. Now we have 
\begin{align*}
    m_{\tau}& =  e(\tau) (-1)^{\dim(C_{\phi})} \langle \Ind_M^G (\pi(\phi_M,\tau_M)) , \IC(\locsys{\tau}{C_{\phi}}) \rangle\\
    &=e(\tau) \langle \Ind_M^G (\pi(\phi_M,\tau_M)) , \stsh{\tau}{C_{\phi}} \rangle \\
    &=e(\tau)  \langle \pi(\phi_{M}, \tau_M) , \xi^{\ast} (\stsh{\tau}{C_{\phi}})\rangle\\
    &= e(\tau)  \left\langle \pi(\phi_{M}, \tau_M) , \sum_{\xi (D) \subseteq C_{\phi},\ \sigma \in \irrepdual{A}_{D}} m(\sigma, \tau) \stsh{\sigma}{D}\right\rangle\\
    &= e(\tau) e(\tau_M) m(\tau_M, \tau)\\
    &= m(\tau_M, \tau).
\end{align*}
Here, the second and the fifth equalities follow from the same reason as \eqref{eq pairing converse}; the third equality is our assumption \eqref{equation: ind res adjunction}; the fourth equality is \eqref{eqn: restriction formula on standards}; and the last equality is \eqref{eqn Kottwitz sign Levi}. This completes the proof of Part (b) and the proposition.
\end{proof}

\iffalse
\section{Perverse sheaves corresponding to Supercuspidal representations}
A representation $\pi\in \Pi(G)$ is supercuspidal if it does not occur in $\Ind_P^G(\sigma)$ for any proper parabolic subgroup $P=MN$ of $G$ and $\sigma \in \Pi(M)$ is inflated to be a representation of $P$. Suppose $\lambda$ is the infinitesimal parameter of the Langlands parameter of $\pi$. Then, we get
\[\langle \Ind_M^G(\sigma), \mathcal{P}(\pi)\rangle = 0 ~~ \forall \text{ proper Levi sugbroups }M, \forall \sigma \in \Pi_{\lambda_M}(M).\] \fi
% \todo[inline]{exposition can be improved by packaging the stuff below in a lemma since it is also used in the proof above, MR can do it once we finalize other things or we can skip it.}
% The same identity also recovers the second of the two conjectures discussed in the introduction. We use the following observation twice, once for $G$ and once for $M$. Let $H$ be a quasisplit connected reductive algebraic group over $F$ equipped with an $\LLC^{\mathfrak{e}}$ satisfying Desideratum \ref{desideratum: Langlands classification}, let $\lambda'$ be an infinitesimal parameter of $H$, and let $\rho=\pi(\phi_\rho,\tau_\rho) \in \Pi^{\pure}_{\lambda'}(H)$ be tempered. Then $\phi_\rho$ is tempered, so $C_{\phi_\rho}$ is open in $V_{\lambda'}$ by \cite[Proposition 3.2]{cunningham2024whittakernormalizationpadicabvpackets}; furthermore this open orbit is the unique one in the prehomogenous vector space $V_{\lambda'}$. In particular $\phi_\rho$ is the unique tempered $L$-parameter of $H$ with infinitesimal parameter $\lambda'$. From \eqref{equation: standard sheaf as a sum of perverse sheaf and boundary terms}, for every $H_{\lambda'}$-orbit $D \subseteq V_{\lambda'}$ and every $\sigma \in \irrepdual{A}_D$,
% \begin{equation}\label{equation: tempered pairing on standards}
%     \big\langle \rho,\stsh{\sigma}{D}\big\rangle=\begin{cases} e(\rho) & (D,\sigma)=(C_{\phi_\rho},\tau_\rho),\\ 0 & \text{otherwise},\end{cases}
% \end{equation}
% Indeed, $\stsh{\sigma}{D}$ is supported on $\ol{D}$, so its class is a combination of classes of simple perverse sheaves supported on $\ol{D}$; if $D\neq C_{\phi_\rho}$ these do not include $\mathcal{IC}(\locsys{\tau_\rho}{C_{\phi_\rho}})$, because $C_{\phi_\rho}$ is open and dense, and the pairing vanishes by \eqref{eqn:pairing}. If $D=C_{\phi_\rho}$, then the extension by zero and the perverse extension of $\locsys{\sigma}{C_{\phi_\rho}}$ agree up to the sign $(-1)^{\dim C_{\phi_\rho}}$ modulo classes supported on smaller orbits, and that sign cancels the sign $(-1)^{d(\rho)}$ of \eqref{eqn:pairing}. Thus the pairing formula \eqref{equation: pairing on standards} of the Kazhdan--Lusztig hypothesis is automatic on tempered representations. The same argument is used in the proof of Proposition \ref{theorem: ind res adjunction implies klh} once we reduce to the tempered case using identity \eqref{equation: ind res adjunction}.

%  {\begin{prop}\label{corollary: converse tempered}
% Let $G$, $M$ and $\lambda$ be as in Proposition \ref{theorem: ind res adjunction implies klh}, and assume the identity \eqref{equation: ind res adjunction}. Then Desideratum \ref{desideratum: induction of tempered representations} holds: for every tempered $L$-parameter $\phi_M$ of $M$ with infinitesimal parameter $\lambda_M$ and every $\tau_M \in \irrepdual{A}_{\phi_M}$, writing $\phi=\xi\circ\phi_M$,
% \[\Ind_{M}^{G} \big(\pi(\phi_M,\tau_M)\big)=\bigoplus_{\tau \in \irrepdual{A}_{\phi}} m(\tau_M, \tau)\, \pi(\phi, \tau).\]
% \end{prop}
% \begin{proof}
% Note that $\phi$ is a tempered $L$-parameter of $G$ with infinitesimal parameter $\lambda$.

% Parabolic induction of a tempered representation is a semisimple sum of tempered representations, so by \eqref{equation: induction definition} we may write
% \[\Ind_M^G\big(\pi(\phi_M,\tau_M)\big)=\sum_{\rho} n_\rho\,[\rho], \qquad n_\rho \in \Z_{\geq 0},\]
% the sum running over the tempered $\rho \in \Pi_\lambda^{\pure}(G)$. By the uniqueness just noted, every such $\rho$ has $\LLC(\rho)=\phi$, since $\phi$ is the unique tempered $L$-parameter of $G$ with infinitesimal parameter $\lambda$; that is, the constituents already lie in the expected $L$-packet, and only their multiplicities are at issue. Writing $n_{\tau_0}:=n_{\pi(\phi,\tau_0)}$ for $\tau_0 \in \irrepdual{A}_\phi$, it therefore suffices to show that $n_{\tau_0}=m(\tau_M,\tau_0)$ for all $\tau_0$.

% Fix $\tau_0 \in \irrepdual{A}_\phi$ and pair both sides with the standard sheaf $\stsh{\tau_0}{C_\phi}$ on the open orbit. On the representation side, \eqref{equation: tempered pairing on standards} gives
% \[\big\langle \Ind_M^G(\pi(\phi_M,\tau_M)),\stsh{\tau_0}{C_\phi}\big\rangle=n_{\tau_0}\,e\big(\kappa_G(\tau_0)\big).\]
% On the geometric side, the restriction formula \eqref{eqn: restriction formula on standards} gives
% \[\xi^*\big(\stsh{\tau_0}{C_\phi}\big)=\sum_{d\xi(D)\subseteq C_\phi,~\sigma \in \irrepdual{A}_D} m(\sigma,\tau_0)\,\stsh{\sigma}{D},\]
% and $\pi(\phi_M,\tau_M)$ is a tempered member of $\Pi_{\lambda_M}^{\pure}(M)$, so \eqref{equation: tempered pairing on standards}, applied this time to $M$, kills every term except the one with $(D,\sigma)=(C_{\phi_M},\tau_M)$ {, a term which does occur because $C_\phi$ is the saturation of $C_{\phi_M}$}. Hence
% \[\big\langle \pi(\phi_M,\tau_M),\xi^*(\stsh{\tau_0}{C_\phi})\big\rangle=m(\tau_M,\tau_0)\,e\big(\kappa_M(\tau_M)\big).\]
% The identity \eqref{equation: ind res adjunction} equates the two displays:
% \[n_{\tau_0}\,e\big(\kappa_G(\tau_0)\big)=m(\tau_M,\tau_0)\,e\big(\kappa_M(\tau_M)\big).\]
% If $m(\tau_M,\tau_0)=0$ this forces $n_{\tau_0}=0$. If $m(\tau_M,\tau_0)\neq 0$ then $\kappa_G(\tau_0)=H_\xi(\kappa_M(\tau_M))$ by Lemma \ref{lemma: relating galois colomology class of irreps of component groups}, so the two Kottwitz signs agree by \eqref{eqn Kottwitz sign Levi} and $n_{\tau_0}=m(\tau_M,\tau_0)$. As the induction is semisimple, this equality of classes in $K\Pi_\lambda^{\pure}(G)$ is an isomorphism of representations, which is \eqref{eq decomp of tempered induction}.
% \end{proof}}

%\todo[inline]{this section seems "too trivial" so we can decide whether its useful; i wonder if we know of a small example with levi subgroups where restriction mixes pure inner forms. The one I have is a non levi endoscopic group\\ \textbf{[disposition:} kept, with an opening remark explaining that $H_\xi$ is bijective here so no mixing occurs, and a closing remark checking the transpose relation. A mixing example needs a non-injective $H_\xi$; not constructed.\textbf{]}}


\section{Supercuspidal sheaves}\label{section: supercuspidal sheaves}
In this section, we use the geometric description of parabolic induction \eqref{equation: ind res adjunction} to 
describe the simple perverse sheaves corresponding to supercuspidal representations under local Langlands correspondence \eqref{eqn: vogan langlands correspondence}.


\begin{defn}\label{def sc sheaf}
A simple perverse sheaf $\mathcal{P} \in \Per_{H_\lambda}(V_\lambda)$ is called  \emph{supercuspidal} if for every proper Levi subgroup $M$ for which $\lambda$ factors through the embedding $\xi: {}^L M \hookrightarrow {}^L G$ as $\lambda_M$, we have
\begin{equation}\label{equation: supercuspidal sheaf condition}
    \xi^*(\mathcal{P})=0 \quad \text{in } K\Per_{H_{\lambda_M}}(V_{\lambda_M}).
\end{equation}
\end{defn}

\begin{prop}
    Keep the assumptions in Proposition \ref{proposition: pure parabolic induction and restriction} for $G$. A representation $\pi \in \Pi_{\lambda}^{\pure}(G)$ is supercuspidal if and only if the corresponding simple perverse sheaf $\mathcal{P}(\pi)\in \Per_{H_\lambda}(V_\lambda)$ is supercuspidal.
\end{prop}
\begin{proof}
    We prove that $\pi$ is not supercuspidal if and only if $\mathcal{P}(\pi)$ is not supercuspidal.

    Suppose that $\pi$ is not supercuspidal, then there exists a proper Levi subgroup $M$ and $\sigma \in \Pi^{\pure}(M)$ such that $\pi$ is a subquotient of $\Ind_M^G(\sigma)$.  Namely, 
    \[0 \neq \langle \Ind_M^G(\sigma), \mathcal{P}(\pi) \rangle.\]
    Write $\lambda_M$ for the infinitesimal parameter of $\sigma$ and write $\xi: {}^LM \hookrightarrow {}^L G$ for the embedding. Then Desideratum \ref{desideratum: Haines} implies that $({}^L M \hookrightarrow {}^L G) \circ \lambda_M= \lambda$. Now Proposition \ref{proposition: pure parabolic induction and restriction} gives
    \[  0 \neq \langle \Ind_M^G(\sigma), \mathcal{P}(\pi) \rangle=  \langle \sigma, \xi^{\ast}(\mathcal{P}(\pi)) \rangle. \]
    Therefore, we must have $ \xi^{\ast}(\mathcal{P}(\pi))\neq 0$ and $\mathcal{P}(\pi)$ is not supercuspidal.

    Conversely, suppose that $\mathcal{P}(\pi)$ is not supercuspidal. Take a Levi subgroup $M$ such that $\xi^{\ast} (\mathcal{P}(\pi)) \neq 0$ where $\xi:{}^L M \hookrightarrow {}^L G$. Then there exists a $\sigma \in \Pi_{\lambda_M}^{\pure}(M)$ such that $\langle \sigma ,\xi^{\ast} (\mathcal{P}(\pi))  \rangle \neq 0 $. Then again by Proposition \ref{proposition: pure parabolic induction and restriction}, we obtain
    \[ 0 \neq \langle \sigma ,\xi^{\ast} (\mathcal{P}(\pi))\rangle = \langle \Ind_M^G \sigma , \mathcal{P}(\pi) \rangle. \]
    This implies that $\pi$ appears in the induction $\Ind_M^G \sigma$, and hence $\pi$ is not supercuspidal. This completes the proof of the proposition.
\end{proof}

\begin{remark} If $\lambda$ does not factor through the $L$-group of any proper Levi subgroups of $G$, then every simple perverse sheaf in $\Per_{H_\lambda}(V_{\lambda})^{\mathrm{simple}}_{\mathrm{/iso}}$ is supercuspidal by Definition \ref{def sc sheaf}. In this case, all of the representations in $\Pi_{\lambda}^{\pure}(G)$ (which turns out to be an $L$-packet with $L$-parameter $\lambda$) are supercuspidal by Desideratum \ref{desideratum: Haines} as in the proof above.   
\end{remark}

% \section{Supercuspidal sheaves}\label{section: supercuspidal sheaves}
% In this section we use the pairing \eqref{eqn:pairing} and the identity $\Ind_M^G={}^t\xi^*$ of Theorem \ref{theorem: main theorem in the introduction} to characterize the simple perverse sheaves corresponding to supercuspidal representations.

% Let $\delta \in H^1(F,G)$. Recall that $\pi \in \Pi(G_\delta)$ is supercuspidal if $\pi$ is not isomorphic to a subquotient of $\Ind_P^{G_\delta}(\sigma)$ for any proper parabolic subgroup $P$ of $G_\delta$ with Levi component $M$ and any $\sigma \in \Pi(M)$. 

% \todo[inline]{I don't want to define supercuspidal for virtual rep. A definition I have seen from Aubert talk is all Jacquet modules are zero, and I am not sure if it is equivalent to the one you wrote below.


% In the notation of \eqref{equation: induction definition}, we write that $\pi \in K\Pi_\lambda(G_\delta)$ is supercuspidal if and only if $\pi$ is not a summand of $\Ind_{M_\gamma}^{G_\delta}(\sigma) \in K\Pi^\pure(G)$ for any proper standard Levi subgroup $M_\gamma$ of $G_\delta$ ( $H_\xi(\gamma)=\delta$) and for any $\sigma \in K\Pi(M)$.
% }

% Now, let us invoke $\LLC^{\mathfrak{e}}$ and let $\lambda$ be an infinitesimal parameter of $G$. Suppose $\lambda$ does not factor through any Levi subgroup of $G$. Then, $\lambda$ must be irreducible (i.e., it should correspond to an irreducible representation $\lambda^0:W_F \to \dualgroup{G}(\CC)$ of $W_F$ with image in $\dualgroup{G}(\CC)$.) If $\phi$ is any Langlands parameter so that $\phi(w,d_w)=\lambda(w)$ for all $w\in W_F$, then it forces $\phi\vert_{\SL_2(\CC)}=1$ and $\phi=\lambda$. Thus, there is only one $L$-parameter $\phi$ with infinitesimal parameter $\lambda$ and $\Pi^\pure_\lambda(G)=\Pi_\phi^\pure(G)$ is an $L$-packet. We impose the following desideratum.
% \begin{desideratum}\label{desideratum: supercuspidal}
%     If an $L$-parameter $\phi$ corresponds to an irreducible representation of $W_F$ with trivial restriction to $\SL_2(\CC)$, then every element of $\Pi_\phi^\pure(G)$ is supercuspidal.
% \end{desideratum}


% Now for a Levi subgroup $M$ of $G$ (with appropriate subscripts if one is specifying a pure inner form), write
% \[\xi_M: \Lgroup{M} \hookrightarrow \Lgroup{G}.\]
% Thus, for an infinitesimal parameter $\lambda$ of $G$ that factors through an infinitesimal parameter $\lambda_M$ of $M$, we write $\lambda=\xi_M \circ \lambda_M$ (up to $\dualgroup{G}(\CC)$-conjugacy). This gives us the corresponding notions of $V_{\lambda_M}$ and $H_{\lambda_M}$, and $\xi^*_M: K\Per_{H_\lambda}(V_\lambda) \to K\Per_{H_{\lambda_M}}(V_{\lambda_M})$. {\color{red}Crucially, if $\lambda$ does not factor through a given Levi subgroup $L$ of $G$, write $\Pi_{\lambda_L}^\pure(M)=\emptyset$; likewise $V_{\lambda_L}=\emptyset$ so that the pullback $\xi^*_L \equiv 0$, where $0$ denotes the trivial sheaf on the empty set.} Now, assuming desideratum \ref{desideratum: supercuspidal}, the notion of supercuspidality can be formulated in terms of the pairing \eqref{eqn:pairing}:
% \begin{align}\label{equation: supercuspidal}
%     \pi \in K\Pi_\lambda^\pure(G) \text{ is supercuspidal} &\iff \langle \Ind_M^G(\sigma), \mathcal{P}(\pi)\rangle = 0 ~~\text{ for all } \sigma \in K\Pi_{\lambda_M}^\pure(M), \notag \\ &\text{ as } M \text{ varies over all proper Levi subgroups }\text{ of }G. 
% \end{align}  

% We may now formulate the notion of a supercuspidal sheaf.
% \begin{defn}\label{defn: supercuspidal sheaf}
% A perverse sheaf $\mathcal{P} \in K\Per_{H_\lambda}(V_\lambda)$ is called  \emph{supercuspidal} if
% \begin{equation}\label{equation: supercuspidal sheaf condition}
%     \xi_{M}^*(\mathcal{P})=0 \quad \text{in } K\Per_{H_{\lambda_M}}(V_{\lambda_M})
% \end{equation}
% for every proper standard Levi subgroup $M$ of $G$. 
% \end{defn}

% Recall the perverse sheaf $\mathcal{P}(\pi) \in K\Per_{H_\lambda}(V_\lambda)$ associated to $\pi \in K\Pi_\lambda^\pure(G)$ via the geometric perspective on $\LLC^{\mathfrak{e}}$ from Section \ref{Vogan--Langlands correspondence}. We now justify the term supercuspidal for the sheaf.
% \begin{prop}
%     Assume the hypothesis of Proposition \ref{proposition: pure parabolic induction and restriction}. Additionally, assume desideratum \ref{desideratum: supercuspidal}. An irreducible representation $\pi \in \Pi_\lambda^\pure(G)$ is supercuspidal if and only if $\mathcal{P}(\pi)\in K\Per_{H_\lambda}(V_\lambda)$ is supercuspidal. 
% \end{prop}
% \begin{proof}
%     Suppose $\pi \in \Pi_\lambda(G)$ is a supercuspidal representation. From \eqref{equation: supercuspidal} which crucially uses desideratum \ref{desideratum: supercuspidal}, this is equivalent to $\langle \Ind_M^G(\sigma), \mathcal{P}(\pi)\rangle = 0$ for all $\sigma \in K\Pi_{\lambda_M}^\pure(M)$  as $ M$  varies over all proper Levi subgroups $G$. Using the identity $\Ind_M^G={}^t\xi^*_M$ from the conclusion of Proposition \ref{proposition: pure parabolic induction and restriction}, we get 
%     \begin{align*}
%         0 &= \langle \Ind_M^G(\sigma), \mathcal{P}(\pi)\rangle\\
%         &=\langle \sigma, \xi^*_M\mathcal{P}(\pi)\rangle
%     \end{align*}
%     for all $\sigma \in \Pi_{\lambda_M}^\pure(M)$ as $M$ varies over all proper Levi subgroups $M$ of $G$. Since the pairing is non-degenerate, this is equivalent to $\xi^*_M(\mathcal{P}(\pi))=0$ for all proper Levi subgroups $M$ of $G$. In other words, $\pi$ is a supercuspidal representation if and only if $\mathcal{P}(\pi)$ is a supercuspidal perverse sheaf. 
% \end{proof}
% {\color{red}
% %------------------------
% \iffalse
% Every proper parabolic subgroup of $G_\delta$ is $G_\delta(F)$-conjugate to a standard one, whose Levi component is $M_\gamma$ for a proper standard Levi subgroup $M$ of $G$ and some $\gamma \in H^1(F,M)$ with $H_\xi(\gamma)=\delta$ (Section \ref{subsection: setting}), and conjugation does not change the class of the induced representation in $K\Pi(G_\delta)$ (Section \ref{subsec: parabolic induction}). Moreover, if $\sigma \in \Pi(M_\gamma)$ with $H_\xi(\gamma) \neq \delta$, then $\Ind_M^G(\sigma)$ is a representation of a pure inner form of $G$ different from $G_\delta$, so it does not contain $\pi$. Hence, in the notation of \eqref{equation: induction definition},
% \begin{equation}\label{equation: supercuspidal via pure induction}
%     \pi \text{ is supercuspidal} \quad \Longleftrightarrow \quad \pi \text{ is not a subquotient of } \Ind_M^G(\sigma) \text{ for all proper standard Levi } M \text{ of } G \text{ and all } \sigma \in \Pi^\pure(M).
% \end{equation}

% Now fix an infinitesimal parameter $\lambda$ of $G$. For a proper standard Levi subgroup $M$ of $G$, let $\xi_M: {}^LM \hookrightarrow {}^LG$ be the $L$-embedding \eqref{eq Levi L embedding}, and let $\Lambda_M(\lambda)$ be the set of infinitesimal parameters $\lambda_M$ of $M$ such that $\xi_M \circ \lambda_M$ is $\dualgroup{G}$-conjugate to $\lambda$. Since $\xi_M$ is well defined up to $\dualgroup{G}$-conjugation, for each $\lambda_M \in \Lambda_M(\lambda)$ we may and do choose $\xi_M$ so that $\xi_M\circ\lambda_M=\lambda$, as in \eqref{equation: setting inf par factors}. We write
% \[ \xi_{M,\lambda_M}^*: K\Per_{H_\lambda}(V_\lambda) \to K\Per_{H_{\lambda_M}}(V_{\lambda_M}) \]
% for the map \eqref{equation: restriction linear map perverse sheaves}, with the dependence on $M$ and $\lambda_M$ made explicit.

% \begin{defn}\label{defn: supercuspidal sheaf}
% A simple perverse sheaf $\mathcal{P} \in \Per_{H_\lambda}(V_\lambda)^{\mathrm{simple}}_{/\mathrm{iso}}$ is called a \emph{supercuspidal sheaf} if
% \begin{equation}\label{equation: supercuspidal sheaf condition}
%     \xi_{M,\lambda_M}^*(\mathcal{P})=0 \quad \text{in } K\Per_{H_{\lambda_M}}(V_{\lambda_M})
% \end{equation}
% for every proper standard Levi subgroup $M$ of $G$ and every $\lambda_M \in \Lambda_M(\lambda)$. Equivalently, since the pairing \eqref{eqn:pairing} for $M$ is perfect, $\mathcal{P}$ is a supercuspidal sheaf if and only if
% \begin{equation}\label{equation: supercuspidal sheaf condition pairing}
%     \langle \sigma, \xi_{M,\lambda_M}^*(\mathcal{P}) \rangle = 0
% \end{equation}
% for all such $M$ and $\lambda_M$ and all $\sigma \in K\Pi^\pure_{\lambda_M}(M)$.
% \end{defn}

% \begin{remark}
% The condition \eqref{equation: supercuspidal sheaf condition} concerns the class of $\xi^*_{M,\lambda_M}\mathcal{P}$ in the Grothendieck group, not the complex $\xi^*_{M,\lambda_M}\mathcal{P} \in D_{H_{\lambda_M}}(V_{\lambda_M})$ itself. It is a purely geometric condition on $V_\lambda$: writing $\mathcal{P}$ in the basis of standard sheaves \eqref{equation: gr grp perverse sheaves} using the change of basis matrix $c_\lambda$ of Section \ref{sec: introduction}, the class $\xi^*_{M,\lambda_M}(\mathcal{P})$ is computed by the formula \eqref{eqn: restriction formula on standards} from the orbits $D \subseteq V_{\lambda_M}$ and the branching multiplicities $m(\sigma,\tau)$ for $A_D \to A_C$.
% \end{remark}

% \begin{prop}\label{prop: supercuspidal sheaf}
% Assume the hypotheses of Theorem \ref{theorem: main theorem in the introduction}, and assume that its equivalent conditions \textup{(a)} and \textup{(b)} hold. Let $\pi \in \Pi^\pure_\lambda(G)$. Then $\pi$ is supercuspidal if and only if $\mathcal{P}(\pi)$ is a supercuspidal sheaf.
% \end{prop}
% \begin{proof}
% Let $M$ be a proper standard Levi subgroup of $G$ and let $\sigma \in \Pi^\pure(M)$ have infinitesimal parameter $\lambda_M$. By Desideratum \ref{desideratum: Haines}, every irreducible subquotient of $\Ind_M^G(\sigma)$ lies in $\Pi^\pure_{\xi_M\circ\lambda_M}(G)$. Hence, if $\lambda_M \notin \Lambda_M(\lambda)$, then $\pi$ is not a subquotient of $\Ind_M^G(\sigma)$. If $\lambda_M \in \Lambda_M(\lambda)$, then by \eqref{eqn:pairing} and Theorem \ref{theorem: main theorem in the introduction}(b), the multiplicity of $\pi$ in $\Ind_M^G(\sigma)$ is
% \[ (-1)^{d(\pi)} e(\pi) \langle \Ind_M^G(\sigma), \mathcal{P}(\pi) \rangle = (-1)^{d(\pi)} e(\pi) \langle \sigma, \xi_{M,\lambda_M}^*(\mathcal{P}(\pi)) \rangle. \]
% Combining this with \eqref{equation: supercuspidal via pure induction}, $\pi$ is supercuspidal if and only if $\langle \sigma, \xi_{M,\lambda_M}^*(\mathcal{P}(\pi)) \rangle=0$ for all proper standard Levi subgroups $M$, all $\lambda_M \in \Lambda_M(\lambda)$ and all $\sigma \in \Pi^\pure_{\lambda_M}(M)$. Since $\Pi^\pure_{\lambda_M}(M)$ spans $K\Pi^\pure_{\lambda_M}(M)$, this is the condition \eqref{equation: supercuspidal sheaf condition pairing}, which is equivalent to \eqref{equation: supercuspidal sheaf condition} by the perfectness of the pairing. This proves the proposition.
% \end{proof}

% \begin{remark}
% In particular, under the assumptions of Proposition \ref{prop: supercuspidal sheaf}, whether $\pi \in \Pi^\pure_\lambda(G)$ is supercuspidal can be read off from the geometry of the Vogan varieties $V_\lambda$ and $V_{\lambda_M}$ alone, without reference to the representation $\pi$ itself.
% \end{remark}

% \fi}


\bibliographystyle{amsalpha} 
\bibliography{references} 
\end{document}


% \section{Example} {\label{section: example}}
% {We illustrate Theorem \ref{theorem: main theorem in the introduction} in the smallest interesting case. In this example the map $H_\xi$ from \eqref{equation: definition of h xi} is bijective, so the phenomenon of $\xi^*$ mixing pure inner forms, discussed after Theorem \ref{theorem: main theorem in the introduction}, does not occur; the example is nevertheless a complete and explicit verification of the transpose relation.}

% Let $G$ be split $\SO_7$, so $\dualgroup{G}=\Sp_6$. Let $M=\GL_2\times \SO_3$ be a Levi subgroup of $\SO_7$, so $\dualgroup{M}=\GL_2\times \Sp_2$.  {The local Langlands correspondence} along with Desiderata \ref{desideratum: Langlands classification} {,} \ref{desideratum: induction of tempered representations}  {and \ref{desideratum: isomorphism}} is known for general linear groups due to \cite{HT01} and for classical groups due to \cite{Art13}.  {Hypothesis} \ref{hypothesis: KLH} is known for these groups due to \cite{Sol25}.

% We note that $H^1(F,G)=\{\pm 1\} \simeq H^1(F,M)$. We also use $\xi: \GL_2(\CC)\times \Sp_2(\CC)\hookrightarrow \Sp_6(\CC)$. Setting $d_w=\left(\begin{smallmatrix}
%     |w|^{1/2}&\\&|w|^{-1/2}
% \end{smallmatrix}\right)$, define $\lambda: W_F\to \Sp_6(\CC)$ by $\xi\circ\lambda_M$, where $\lambda_M:W_F \to \GL_2(\CC)\times \Sp_2(\CC)$ is $\lambda_M(w)=d_w\boxtimes d_w$. 

% \begin{lemma}\label{lemma: orbits for G}
%     The Vogan variety $V_\lambda \simeq \A^6$ whose elements can be identified with matrices 
% $X$ of the form $\left(\begin{smallmatrix}
%     b&a&c\\d&f&a\\e&d&b
% \end{smallmatrix}\right)$. The stabilizer group $H_\lambda \simeq \GL_3(\CC)$. The $H_\lambda$-action on $V_\lambda$ is given by $A\cdot X=AXA^{-1}$, which yields $4$ orbits $C_i$ consisting of matrices of rank $i$ for $i=0,1,2,3$.
% \end{lemma}
% \begin{proof}
% The proof follows from definitions. For ease of computation, one can define $\Sp_6$ using the bilinear form given by the matrix $\left(\begin{smallmatrix}
%     &J\\-J
% \end{smallmatrix}\right)$
% where $J=\left(\begin{smallmatrix}
%     &&1\\&1&\\1&&
% \end{smallmatrix}\right)$. Then, $\lambda(w)=\left(\begin{smallmatrix}
%     |w|^{1/2}I_3&\\&|w|^{-1/2}I_3
% \end{smallmatrix} \right)$ and the rest follows easily.\end{proof} 
% For an $H_\lambda$-orbit $C$ in $V_\lambda$, choose a basepoint $x\in C$ to compute the component group $A_C=\pi_0(Z_{H_\lambda}(x))$. For an irreducible representation $\tau \in \irrepdual{A}_C$, recall the Kottwitz sign $e(\kappa_G(\tau))=\pm 1$ from  {\eqref{equation: kottwitz sign}}, where $\kappa_G(\tau) \in H^1(F,G)$ is defined in \eqref{eqn: kappa tau}. With \eqref{eq: simple perverse sheaves classified} in mind, we have the lemma below.
% \begin{lemma}\label{lemma: component groups and local systems G}
%     The following table summarizes the data of component groups and local systems for the orbits from Lemma \ref{lemma: orbits for G}.
%    \begin{center}
% \begin{tabular}{|c|c|c|c|}
% \hline
%    $C$  & $\hat{A}_C$ & Local systems ($e=1$) &Local systems $(e=-1)$  \\
%      \hline
%     $C_0$ & $1$ &$\mathbbm{1}_{C_0}$&\\
%     \hline
%     $C_1$ & $\mathbb{Z}/2\mathbb{Z}$ &$\mathbbm{1}_{C_1}$&$\mathcal{L}_{C_1}$\\
%     \hline
%      $C_2$ & $\mathbb{Z}/2\mathbb{Z}$ &$\mathbbm{1}_{C_2}, \mathcal{L}_{C_2}$&\\
%     \hline
%      $C_3$ & $\Z/2\Z$ &$\mathbbm{1}_{C_3}$&$\mathcal{L}_{C_3}$\\
%      \hline
% \end{tabular}  
% \end{center}
%     \end{lemma}
% \begin{proof}
%     This is a routine computation, but we demonstrate this for the case of the $H_\lambda$-orbit $C_3$ consisting of $3\times 3$-matrices of full rank. Choose the basepoint $x=\left(\begin{smallmatrix}
%         &&1\\&1&\\1&&
%     \end{smallmatrix}\right)$. Then $Z_{H_\lambda}(x)\simeq O(3,J)$, the group of $3\times 3$ orthogonal matrices with respect to $J=x$. It has two connected components corresponding to determinant $1$ and $-1$, so $A_{C_3}=\pi_0(Z_{H_\lambda}(x))\simeq \Z/2\Z$. Then $\irrepdual{A}_{C_3}=\{1,\sgn\}$ which gives us the equivariant local systems $\1_{C_3}$ and $\mathcal{L}_{C_3}$. Next, the map $\{\pm I_{6}\}=Z(\dualgroup{G}(\CC))\to Z_{H_{\lambda}}(x)$ which sends the identity matrix $I_6$ (respectively, $-I_6$) to the connected component of determinant $1$ (respectively, $-1$) gives the  {bijection} $\irrepdual{A}_{C_3} \simeq  {\Hom(\pi_0(Z(\dualgroup{G})^{\Gamma_F}), \CC^\times)}$. Using the map \ref{eqn: kottwitz map}, $\sgn$ corresponds to the nontrivial pure inner form. 
% \end{proof}
% Using definitions, one can similarly obtain the geometric data for $M$. 
% \begin{lemma}
%     The Vogan variety for $M$ is\[V_{\lambda_M}=  \left\{\begin{bmatrix}0&a\\0&0\end{bmatrix} \boxtimes \begin{bmatrix}0&e\\0&0\end{bmatrix}\right\} \simeq \mathbb{A}^2,\] and the stabilizer group $H_{\lambda_M}$ is 
%     \[H_{\lambda_M}= \left\{\begin{bmatrix}t_1&0\\0&t_2\end{bmatrix} \boxtimes \begin{bmatrix}t&0\\0&t^{-1}\end{bmatrix}\right\} \simeq (\BC^{\times})^3.\]
%     The $H_{\lambda_M}$-action on $V_{\lambda_M}$ is given by \[(t_1,t_2,t)\cdot (a,e)=\left(\frac{t_1}{t_2}a,t^2e\right).\]There are $4$ orbits, depending on whether $a$ and $e$ are $0$ or nonzero. We denote the orbits as $C_{00}, C_{10}, C_{01}$, and $C_{11}$, where the notation $C_{ij}$ means the orbit of the point $(a,e)=(i,j)$.
% \end{lemma}

% \begin{lemma}\label{lemma: component groups and local systems levi subgroup}
% As in Lemma \ref{lemma: component groups and local systems G}, we write down the component groups and equivariant local systems in the table below.
%     \begin{center}
% \begin{tabular}{|c|c|c|c|}
% \hline
%    $C$  & $\hat{A}_C$ & Local systems($e_M=1$) &Local Systems($e_M=-1$) \\
%      \hline
%     $C_{00}$ & $1 $ &$\mathbbm{1}_{C_{00}}$&\\
%     \hline
%     $C_{01}$ & $\mathbb{Z}/2\mathbb{Z}$ &$\mathbbm{1}_{C_{01}}$ & $\mathcal{L}_{C_{01}}$\\
%     \hline
%      $C_{10}$ & $1$ &$\mathbbm{1}_{C_{10}}$&\\
%     \hline
%      $C_{11}$ & $\mathbb{Z}/2\mathbb{Z}$ &$\mathbbm{1}_{C_{11}}$ & $\mathcal{L}_{C_{11}}$ \\
%      \hline
% \end{tabular}  
% \end{center}
% \end{lemma}
% \begin{proof}
%     The proof follows from routine computation. See proof of Lemma \ref{lemma: component groups and local systems G} for an instance of such a computation.
% \end{proof}
% The map $\xi$ induces $\xi: H_{\lambda_M} \hookrightarrow H_{\lambda}$ and $d\xi: V_{\lambda_M} \hookrightarrow V_{\lambda}$:
% \begin{equation}\label{equation: L embedding example}
%     (t_1,t_2,t) \mapsto \begin{bmatrix}
%     t_1&&&&&\\&t_2^{-1}&&&&\\&&t&&&\\&&&t^{-1}&&\\&&&&t_2&\\&&&&&t_1^{-1}
% \end{bmatrix}, \hspace{0.5cm}(a,e) \mapsto \begin{bmatrix}
%     0&a&0\\0&0&a\\e&0&0
% \end{bmatrix}
% \end{equation}
% Next, we compute $\mathbb{Z}$-linear map 
% \[\xi^*: K\Per_{H_{\lambda}} (V_{\lambda})\to K\Per_{H_{\lambda_M}}(V_{\lambda_M}). \]
% Let us specify an ordered basis for each Grothendieck group. For $K\Per_{H_{\lambda}} (V_{\lambda})$, we have the basis of constructible standard sheaves (we abusively use the notation of local systems on an orbit to denote sheaves on $V_\lambda$)
% \[\{[\mathbbm{1}_{C_0}], [\mathbbm{1}_{C_1}], [\mathbbm{1}_{C_2}], [\mathcal{L}_{C_2}], [\mathbbm{1}_{C_3}], -[\mathcal{L}_{C_1}], -[\mathcal{L}_{C_3}]\}.\]
%  Note that we have kept the local systems attached to the pure inner form at the end in this ordering. The purpose of using $-1=e(\delta)$ as a coefficient of basis elements associated to the pure inner form $0\neq \delta \in H ^1(F,G)$ of is to make this a biorthogonal dual basis to that of $K\Pi_{\lambda}^\pure(G)$ with respect to the pairing \eqref{equation: pairing on standards}. Next, for $K\Per_{H_{\lambda_M}}(V_{\lambda_M})$, we have the basis of constructible standard sheaves
% \[\{[\mathbbm{1}_{C_{00}}], [\mathbbm{1}_{C_{01}}], [\mathbbm{1}_{C_{10}}], [\mathbbm{1}_{C_{11}}], -[\mathcal{L}_{C_{01}}], -[\mathcal{L}_{C_{11}}] \}\]
% The last two elements of this basis are attached to the pure inner form of the Levi subgroup. Once again, $-1=e(\gamma)$ is a coefficient of basis elements associated to the pure inner form $0\neq \gamma \in H^1(F,M)$ to make this a biorthogonal dual basis to that of $K\Pi_{\lambda}^\pure(M)$ with respect to the pairing \eqref{equation: pairing on standards} for $M$.
% \begin{lemma}\label{lemma: restriction matrix}
%     Using \eqref{eqn: restriction formula on standards}, we can write down $\xi^*$ as the matrix
%   \[\xi^*=\begin{pNiceArray}{ccccc|cc}
%   1 &0 &0 &0 &0&0&0\\
%   0& 1&0&0&0&0&0 \\
%   0&0&1&1&0&0&0\\
%   0&0&0&0&1&0&0 \\
%   0&0&0&0 & 0&1&0\\
%   0&0&0&0  & 0 &0&1
% \end{pNiceArray}     .
% \]
% \end{lemma}
% \begin{proof}
%     Using \eqref{equation: L embedding example}, we can see that $C_{00}=C_0$, $C_{01} \subset C_1$, $C_{10} \subset C_2$, and $C_{11} \subset C_3$. The rest of the computations follow immediately.
% \end{proof}
% Next, we write down the map $\Ind_M^G$. We first begin by identifying all enhanced Langlands parameters associated with orbits associated with $G$ and $M$. Let $S_k$ denote the $k$-dimensional  {irreducible algebraic} representation of $\SL_2(\CC)$. 

% \begin{lemma}\label{lemma: vogan Langlands correspondence for G and M}
%     The following table summarizes the data of the enhanced Langlands parameters associated to split $G$ and its pure inner form, $G'$. Define the Langlands parameters for $G$:
%     \begin{align*}
%     \phi_3&= (\pmb{1} \otimes S_2)^3,\\
%     \phi_2&= |\cdot|^{-1/2} \otimes S_1+ (\pmb{1} \otimes S_2)^2+|\cdot|^{1/2} \otimes S_1,\\
%     \phi_1&= (|\cdot|^{-1/2} \otimes S_1)^2+ \pmb{1} \otimes S_2+(|\cdot|^{1/2} \otimes S_1)^2,\\
%     \phi_0&= (|\cdot|^{-1/2} \otimes S_1)^3+(|\cdot|^{1/2} \otimes S_1)^3.
% \end{align*}
% We have the following correspondence.
%        \begin{center}
% \begin{tabular}{|c|c|c|c|c|}
% \hline
%    $\phi$&$C_\phi$  & $\hat{A}_\phi$ & $\Pi_\lambda(G)$ ($e=1$) & $\Pi_{\lambda}(G')$ $(e=-1)$  \\
%      \hline
%    $\phi_0$ & $C_0$ & $1$ &$\pi_0$&\\
%     \hline
%    $\phi_1$ & $C_1$ & $\mathbb{Z}/2\mathbb{Z}$ &$\pi_1^+$&$\pi_1^-$\\
%     \hline
%    $\phi_2$  & $C_2$ & $\mathbb{Z}/2\mathbb{Z}$ &$\pi_2^+, \pi_2^-$&\\
%     \hline
%    $\phi_3$  & $C_3$ & $\Z/2\Z$ &$\pi_3^+$&$\pi_3^-$\\
%      \hline
% \end{tabular}  
% \end{center}
% \end{lemma}
% Let $\Delta_i^\pm$ denote the standard module associated to $\pi_i^\pm$ {; for example, $\Delta_1^-$} is the standard module of the representation $\pi_1^-\in \Pi_\lambda(G')$. Then, $K\Pi_\lambda^\pure(G)$ has the basis
% \[\{ \Delta_0,\Delta_1^+, \Delta_2^+,\Delta_2^-,\Delta_3^+,\Delta_1^-,\Delta_3^-\}\]
% \begin{lemma}
% Define the following Langlands parameters for quasisplit $M=\GL_2\times \SO_3$:
%     \begin{align*}
%     \phi_{11}&= (\pmb{1} \otimes S_2) \boxtimes (\pmb{1} \otimes S_2),\\
%     \phi_{10}&= (\pmb{1} \otimes S_2 ) \boxtimes (d_w \otimes S_1 ),\\
%     \phi_{01}&= (d_w \otimes S_1 ) \boxtimes (\pmb{1} \otimes S_2 ),\\
%     \phi_{00}&= (d_w \otimes S_1) \boxtimes (d_w \otimes S_1).
% \end{align*}
% Let $M'=\GL_2\times \SO_3'$ denote the pure inner form of $M$. We have the following correspondence.
%     \begin{center}
% \begin{tabular}{|c|c|c|c|c|}
% \hline
%   $\phi$ &$C$  & $\hat{A}_{\phi}$ & $\Pi_{\lambda_M}(M)$($e_M=1$) & $\Pi_{\lambda_M}(M')$($e_M=-1$) \\
%      \hline
%   $\phi_{00}$ & $C_{00}$ & $1 $ &$\pi_{00}$&\\
%     \hline
%   $\phi_{01}$ & $C_{01}$ & $\mathbb{Z}/2\mathbb{Z}$ &$\pi_{01}^+$ & $\pi_{01}^-$\\
%     \hline
%  $\phi_{10}$  &  $C_{10}$ & $1$ &$\pi_{10}$&\\
%     \hline
%  $\phi_{11}$   & $C_{11}$ & $\mathbb{Z}/2\mathbb{Z}$ &$\pi_{11}^+$ & $\pi_{11}^-$ \\
%      \hline
% \end{tabular}  
% \end{center}
% \end{lemma}
%  {Let $\Delta_{ij}^\pm$ be the standard module for $\pi_{ij}^\pm$; for example, $\Delta_{01}^-$ is the standard module for $\pi_{01}^-$.} Then, $K\Pi_{\lambda_M}^\pure(M)$ has the basis
% \[\{\Delta_{00},\Delta_{01}^+, \Delta_{10}, \Delta_{11}^+, \Delta_{01}^-,\Delta_{11}^-\}\]
% \begin{prop}
%     The matrix for $\Ind_M^G$ is given by
%     \[\Ind_M^G=\begin{pNiceArray}{cccc|cc}
%   1 &0 &0 &0 &0&0\\
%   0& 1&0&0&0&0 \\
%   0&0&1&0&0&0\\
%   0&0&1&0&0&0 \\
%   0&0&0 & 1&0&0\\
%   0&0&0  & 0 &1&0\\
%   0&0&0  & 0 &0&1
  
% \end{pNiceArray}     .
% \]
% \end{prop}
%  {\begin{proof}
% This is done by direct calculation of standard modules and their decomposition into irreducible factors in the Grothendieck group.
% \end{proof}}

% {\begin{remark}
% Observe that the matrix of $\Ind_M^G$ just computed is exactly the transpose of the matrix of $\xi^*$ obtained in Lemma \ref{lemma: restriction matrix}, with respect to the two ordered bases fixed above. This is Proposition \ref{proposition: pure parabolic induction and restriction} in this example. Note also that both matrices are block diagonal for the splitting into pure inner forms indicated by the vertical rules; as observed at the beginning of this section, this reflects the fact that $H_\xi:H^1(F,M)\to H^1(F,G)$ is bijective here.
% \end{remark}}
% %\todo[inline]{up to chi-heng\\ \textbf{[disposition:} the final Proposition now has a proof and a closing remark.\textbf{]}}

